% generated from JIRA project LVV
% using template at /Users/womullan/LSSTgit/docsteady/docsteady/templates/tpnoresult.latex.jinja2.
% using docsteady version 2.2.post4+g100285e.d20210519
% Please do not edit -- update information in Jira instead
\documentclass[dm,lsstdraft,STR,toc]{lsstdoc}
\usepackage{geometry}
\usepackage{longtable,booktabs}
\usepackage{enumitem}
\usepackage{arydshln}
\usepackage{attachfile}
\usepackage{array}
\usepackage{dashrule}

\newcolumntype{L}[1]{>{\raggedright\let\newline\\\arraybackslash\hspace{0pt}}p{#1}}

\input meta.tex

\newcommand{\attachmentsUrl}{https://github.com/\gitorg/\lsstDocType-\lsstDocNum/blob/\gitref/attachments}
\providecommand{\tightlist}{
  \setlength{\itemsep}{0pt}\setlength{\parskip}{0pt}}

\setcounter{tocdepth}{4}

\begin{document}

\def\milestoneName{Level 3 System Spread Configuration Integration}
\def\milestoneId{}
\def\product{SIT-COM Integration}

\setDocCompact{true}

\title{LVV-P81: Level 3 System Spread Configuration Integration Test Plan }
\setDocRef{\lsstDocType-\lsstDocNum}
\date{ 2022-07-30 }
\author{ Austin Roberts }

\input{history_and_info.tex}


\setDocAbstract{
This is the test plan for
\textbf{ Level 3 System Spread Configuration Integration},
an LSST milestone pertaining to the Data Management Subsystem.\\
This document is based on content automatically extracted from the Jira test database on \docDate.
The most recent change to the document repository was on \vcsDate.
}


\maketitle

\section{Introduction}
\label{sect:intro}


\subsection{Objectives}
\label{sect:objectives}

 The primary objective of this integration event is to baseline the
interoperability of the available systems located on the third floor of
the Summit prior to the systems populated on the Camera Cart getting
installed on the Telescope Mount Assembly (TMA). The systems populating
the Camera Cart are the Camera Cable Wrap (CCW), Camera Hexapod, Camera
Rotator, and the Commissioning Camera (ComCam). The secondary objective
of this integration event is to check the Global Interlock System (GIS)
and that the connections to the available Interlock Systems on the third
floor are working correctly.\\[2\baselineskip]This integration event is
organized into a series of test cycles. First, we will verify each
system's communication interface with SAL by checking their commands,
events, and telemetry. Then we will verify each of the GIS safety matrix
detections and actuations and the GIS fault conditions. For the
Interlock Systems that are available, we will verify that they receive
the input from the GIS and produce the correct outputs. Finally, we will
start with the M1M3 hardware and the rest of the systems as simulators.
We will run a series of tests in that configuration. Then we will
replace each simulator with the actual hardware that is available one by
one and rerun the series of tests. The end goal of these tests is to run
through 5 observations from the ScriptQueue using ComCam images with
corrections being applied by the Main Telescope Active Optics System
(MTAOS).\\[2\baselineskip]There are a few constraints for this
configuration on the third floor that must be taken into account during
this integration event.\\
The Main Telescope Mount (MTMount) will always be a simulator since the
actual TMA hardware is not available to provide the elevation changes to
the other systems.\\
There are two simulators available. The MTMount simulator and the
\textbf{MTMount hybrid} . The \textbf{MTMount hybrid~} allows for large
(\textgreater{}3deg) movements of the rotator.\\
The used version depends on where the CCW needs to be simulated or not.
CCW needs to be simulated when the rotator is simulated. In this case,
the MTMount simulator is used.\\
The necessary telemetry for the rotator will be provided by both
\textbf{MTMount~} simulators.\\
The M2 Hexapod is already installed on the TMA. Even though it might be
physically available, we are always going to use it as a simulated
component.\\[2\baselineskip]For the M1M3 and the M2 mirror, it is needed
to use the MTMount telemetry when we are using testing the LUT.\\
When only testing the SAL communication only, the inclinometers should
be used for the test case. In future, when the TMA is available and
moving, only inclinometers will be used.\\[2\baselineskip]Few time:
C206, C176, C170, C177\\
Normal time: C170, C177, C173, C172, C174, C206, C175, C176



\subsection{System Overview}
\label{sect:systemoverview}

 The systems involved in this event are as follows:

\begin{itemize}
\tightlist
\item
  MTAOS
\item
  Pointing Component
\item
  ScriptQueue
\item
  MTMount simulator
\item
  M2 Hexapod simulator
\item
  M1M3 hardware
\item
  M2 simulator and hardware
\item
  Camera Hexapod simulator and hardware
\item
  Camera Rotator simulator and hardware
\item
  ComCam hardware
\item
  DM Tools

  \begin{itemize}
  \tightlist
  \item
    Archiver
  \item
    OODS
  \end{itemize}
\end{itemize}

For the GIS the systems involved in this event are as follows:

\begin{itemize}
\tightlist
\item
  GIS Main Controller
\item
  M2 Camera Controller

  \begin{itemize}
  \tightlist
  \item
    M2 Interlock System
  \item
    Camera Rotator Interlock System
  \item
    Camera Hexapod Interlock System
  \end{itemize}
\item
  M1M3 Controller

  \begin{itemize}
  \tightlist
  \item
    M1M3 Interlock System
  \end{itemize}
\end{itemize}


\subsection{Document Overview}
\label{sect:docoverview}

This document was generated from Jira, obtaining the relevant information from the
\href{https://jira.lsstcorp.org/secure/Tests.jspa\#/testPlan/LVV-P81}{LVV-P81}
~Jira Test Plan and related Test Cycles (
\href{https://jira.lsstcorp.org/secure/Tests.jspa\#/testCycle/LVV-C170}{LVV-C170}
\href{https://jira.lsstcorp.org/secure/Tests.jspa\#/testCycle/LVV-C172}{LVV-C172}
\href{https://jira.lsstcorp.org/secure/Tests.jspa\#/testCycle/LVV-C173}{LVV-C173}
\href{https://jira.lsstcorp.org/secure/Tests.jspa\#/testCycle/LVV-C174}{LVV-C174}
\href{https://jira.lsstcorp.org/secure/Tests.jspa\#/testCycle/LVV-C175}{LVV-C175}
\href{https://jira.lsstcorp.org/secure/Tests.jspa\#/testCycle/LVV-C176}{LVV-C176}
\href{https://jira.lsstcorp.org/secure/Tests.jspa\#/testCycle/LVV-C177}{LVV-C177}
\href{https://jira.lsstcorp.org/secure/Tests.jspa\#/testCycle/LVV-C206}{LVV-C206}
).

Section \ref{sect:intro} provides an overview of the test campaign, the system under test (\product{}),
the applicable documentation, and explains how this document is organized.
Section \ref{sect:testplan} provides additional information about the test plan, like for example the configuration
used for this test or related documentation.
Section \ref{sect:personnel} describes the necessary roles and lists the individuals assigned to them.

Section \ref{sect:overview} provides a summary of the test results, including an overview in Table \ref{table:summary},
an overall assessment statement and suggestions for possible improvements.
Section \ref{sect:detailedtestresults} provides detailed results for each step in each test case.

The current status of test plan \href{https://jira.lsstcorp.org/secure/Tests.jspa\#/testPlan/LVV-P81}{LVV-P81} in Jira is \textbf{ Approved }.

\subsection{References}
\label{sect:references}
\renewcommand{\refname}{}
\bibliography{lsst,refs,books,refs_ads,local}


\newpage
\section{Test Plan Details}
\label{sect:testplan}


\subsection{Data Collection}

  Observing is not required for this test campaign.

\subsection{Verification Environment}
\label{sect:hwconf}
  The environment for the execution of this event will be the third floor
of the Summit.

  \subsection{Entry Criteria}
  Due to technical and schedule constraints, the full re-verification of
each of the subsystems integrated into the spread configuration is not
possible.\\
The following functionality of each of the subsystems is required prior
to executing this verification event.

\begin{itemize}
\tightlist
\item
  M2

  \begin{itemize}
  \tightlist
  \item
    Subscribe to MTMount telemetry
  \item
    Apply forces as commanded through SAL
  \item
    Read and apply compensations from the look up table
  \item
    State machine transitions
  \end{itemize}
\item
  M1M3

  \begin{itemize}
  \tightlist
  \item
    Subscribe to MTMount telemetry
  \item
    Apply forces as commanded through SAL
  \item
    Read and apply compensations from the look up table
  \item
    State machine transitions
  \end{itemize}
\item
  Camera Hexapod

  \begin{itemize}
  \tightlist
  \item
    Re-verification of the Camera Hexapod with ComCam installed
  \end{itemize}
\item
  Camera Rotator

  \begin{itemize}
  \tightlist
  \item
    Re-verification of the Camera Rotator with ComCam installed
  \end{itemize}
\item
  ComCam

  \begin{itemize}
  \tightlist
  \item
    Image playback functionality on the Summit
  \item
    Filter change
  \item
    Shutter open and close
  \item
    Simulated image ingested via DAQ for MTAOS integration
  \item
    CCS integrated with SAL on Summit
  \end{itemize}
\item
  MTAOS
\item
  Pointing Component
\item
  ScriptQueue
\end{itemize}

  \subsection{Exit Criteria}
  In order for this event to be considered complete, the following
criteria must be met:

\begin{itemize}
\tightlist
\item
  Raw test data, events, and telemetry have been saved for all systems
  included in the tests.
\item
  All test data has been analyzed and post-processed.
\item
  All test steps have been statused in the Jira Test Cases within this
  Test Plan and actual results populated as required.
\item
  A summary of the results of the test campaign has been captured in the
  Overall Assessment and Recommended Improvements fields of this Test
  Plan
\item
  A link to the verification artifacts used to produce the summary of
  results has been populated in the Verification Artifacts field of this
  Test Plan
\item
  Any failures have been captured in the
  \href{https://jira.lsstcorp.org/projects/FRACAS/issues/}{FRACAS}
  project
\item
  Determine the extent to which all the components operated coherently
  as one system. This means that we can send a command with the
  ScriptQueue to slew, track and observe including the following
  components:

  \begin{itemize}
  \tightlist
  \item
    M2 cell + surrogate on the cart
  \item
    M1M3 cell + surrogate on cart
  \item
    Camera Hexapod
  \item
    Camera Rotator with restrictive motion due to the CCW not being
    available
  \item
    ComCam
  \item
    DM tools
  \end{itemize}
\end{itemize}


\subsection{Related Documentation}

Docushare collection where additional relevant documentation can be found:

\begin{itemize}
\item For the verification artifacts see the executions of the test cases.
\end{itemize}



\subsection{PMCS Activity}

Primavera milestones related to the test campaign:
\begin{itemize}
\item None
\end{itemize}


\newpage
\section{Personnel}
\label{sect:personnel}

The personnel involved in the test campaign is shown in the following table.

{\small
\begin{longtable}{p{3cm}p{3cm}p{3cm}p{6cm}}
\hline
\multicolumn{2}{r}{T. Plan \href{https://jira.lsstcorp.org/secure/Tests.jspa\#/testPlan/LVV-P81}{LVV-P81} owner:} &
\multicolumn{2}{l}{\textbf{ Austin Roberts } }\\\hline
\multicolumn{2}{r}{T. Cycle \href{https://jira.lsstcorp.org/secure/Tests.jspa\#/testCycle/LVV-C170}{LVV-C170} owner:} &
\multicolumn{2}{l}{\textbf{
Holger Drass }
} \\\hline
\textbf{Test Cases} & \textbf{Assigned to} & \textbf{Executed by} & \textbf{Additional Test Personnel} \\ \hline
\href{https://jira.lsstcorp.org/secure/Tests.jspa#/testCase/LVV-T1802}{LVV-T1802}
& {\small Holger Drass } & {\small Holger Drass } &
\begin{minipage}[]{6cm}
\smallskip
{\small (1) Software Engineer\\
(1) Hardware Engineer }
\medskip
\end{minipage}
\\ \hline
\href{https://jira.lsstcorp.org/secure/Tests.jspa#/testCase/LVV-T1615}{LVV-T1615}
& {\small Austin Roberts } & {\small Bruno Quint } &
\begin{minipage}[]{6cm}
\smallskip
{\small (1) Mechanical Engineer\\
(1) Electrical Engineer\\
(1) Software Engineer }
\medskip
\end{minipage}
\\ \hline
\href{https://jira.lsstcorp.org/secure/Tests.jspa#/testCase/LVV-T1600}{LVV-T1600}
& {\small Holger Drass } & {\small  } &
\begin{minipage}[]{6cm}
\smallskip
{\small (1) Software Engineer\\
(1) Hardware Engineer }
\medskip
\end{minipage}
\\ \hline
\href{https://jira.lsstcorp.org/secure/Tests.jspa#/testCase/LVV-T1602}{LVV-T1602}
& {\small Austin Roberts } & {\small  } &
\begin{minipage}[]{6cm}
\smallskip
{\small Integration Specialist\\
Systems Engineer }
\medskip
\end{minipage}
\\ \hline
\href{https://jira.lsstcorp.org/secure/Tests.jspa#/testCase/LVV-T2222}{LVV-T2222}
& {\small Austin Roberts } & {\small  } &
\begin{minipage}[]{6cm}
\smallskip
{\small (1) Integration Specialist\\
(1) Systems Engineer }
\medskip
\end{minipage}
\\ \hline
\href{https://jira.lsstcorp.org/secure/Tests.jspa#/testCase/LVV-T2221}{LVV-T2221}
& {\small Holger Drass } & {\small  } &
\begin{minipage}[]{6cm}
\smallskip
{\small (1) Integration Specialist\\
(1) Systems Engineer }
\medskip
\end{minipage}
\\ \hline
\href{https://jira.lsstcorp.org/secure/Tests.jspa#/testCase/LVV-T2232}{LVV-T2232}
& {\small Petr Kubanek } & {\small Bruno Quint } &
\begin{minipage}[]{6cm}
\smallskip
{\small (1) Mechanical Engineer\\
(1) Electrical Engineer\\
(1) Software Engineer }
\medskip
\end{minipage}
\\ \hline
\href{https://jira.lsstcorp.org/secure/Tests.jspa#/testCase/LVV-T2219}{LVV-T2219}
& {\small Austin Roberts } & {\small Bruno Quint } &
\begin{minipage}[]{6cm}
\smallskip
{\small (1) Mechanical Engineer\\
(1) Electrical Engineer\\
(1) Software Engineer\\
(1) Commissioning Scientist }
\medskip
\end{minipage}
\\ \hline
\multicolumn{2}{r}{T. Cycle \href{https://jira.lsstcorp.org/secure/Tests.jspa\#/testCycle/LVV-C172}{LVV-C172} owner:} &
\multicolumn{2}{l}{\textbf{
Holger Drass }
} \\\hline
\textbf{Test Cases} & \textbf{Assigned to} & \textbf{Executed by} & \textbf{Additional Test Personnel} \\ \hline
\href{https://jira.lsstcorp.org/secure/Tests.jspa#/testCase/LVV-T2193}{LVV-T2193}
& {\small Bo Xin [X] } & {\small  } &
\begin{minipage}[]{6cm}
\smallskip
{\small  }
\medskip
\end{minipage}
\\ \hline
\href{https://jira.lsstcorp.org/secure/Tests.jspa#/testCase/LVV-T2242}{LVV-T2242}
& {\small Bo Xin [X] } & {\small  } &
\begin{minipage}[]{6cm}
\smallskip
{\small Integration Specialist\\
Systems Engineer }
\medskip
\end{minipage}
\\ \hline
\href{https://jira.lsstcorp.org/secure/Tests.jspa#/testCase/LVV-T2213}{LVV-T2213}
& {\small Austin Roberts } & {\small  } &
\begin{minipage}[]{6cm}
\smallskip
{\small Integration specialist\\
Systems Engineer }
\medskip
\end{minipage}
\\ \hline
\href{https://jira.lsstcorp.org/secure/Tests.jspa#/testCase/LVV-T2190}{LVV-T2190}
& {\small Bo Xin [X] } & {\small  } &
\begin{minipage}[]{6cm}
\smallskip
{\small Systems Integration Specialist\\
Systems Verification Engineer }
\medskip
\end{minipage}
\\ \hline
\href{https://jira.lsstcorp.org/secure/Tests.jspa#/testCase/LVV-T2241}{LVV-T2241}
& {\small Bo Xin [X] } & {\small  } &
\begin{minipage}[]{6cm}
\smallskip
{\small Integration Specialist\\
Systems Engineer }
\medskip
\end{minipage}
\\ \hline
\href{https://jira.lsstcorp.org/secure/Tests.jspa#/testCase/LVV-T2214}{LVV-T2214}
& {\small Austin Roberts } & {\small  } &
\begin{minipage}[]{6cm}
\smallskip
{\small Systems Engineer\\
Integration specialist }
\medskip
\end{minipage}
\\ \hline
\href{https://jira.lsstcorp.org/secure/Tests.jspa#/testCase/LVV-T2215}{LVV-T2215}
& {\small Austin Roberts } & {\small  } &
\begin{minipage}[]{6cm}
\smallskip
{\small Systems Verification Engineer\\
System Integration Specialist }
\medskip
\end{minipage}
\\ \hline
\href{https://jira.lsstcorp.org/secure/Tests.jspa#/testCase/LVV-T2216}{LVV-T2216}
& {\small Austin Roberts } & {\small  } &
\begin{minipage}[]{6cm}
\smallskip
{\small (1) Integration Specialist\\
(1) Systems Engineer }
\medskip
\end{minipage}
\\ \hline
\multicolumn{2}{r}{T. Cycle \href{https://jira.lsstcorp.org/secure/Tests.jspa\#/testCycle/LVV-C173}{LVV-C173} owner:} &
\multicolumn{2}{l}{\textbf{
Holger Drass }
} \\\hline
\textbf{Test Cases} & \textbf{Assigned to} & \textbf{Executed by} & \textbf{Additional Test Personnel} \\ \hline
\href{https://jira.lsstcorp.org/secure/Tests.jspa#/testCase/LVV-T2190}{LVV-T2190}
& {\small Holger Drass } & {\small  } &
\begin{minipage}[]{6cm}
\smallskip
{\small Systems Integration Specialist\\
Systems Verification Engineer }
\medskip
\end{minipage}
\\ \hline
\href{https://jira.lsstcorp.org/secure/Tests.jspa#/testCase/LVV-T2193}{LVV-T2193}
& {\small Holger Drass } & {\small  } &
\begin{minipage}[]{6cm}
\smallskip
{\small  }
\medskip
\end{minipage}
\\ \hline
\href{https://jira.lsstcorp.org/secure/Tests.jspa#/testCase/LVV-T2213}{LVV-T2213}
& {\small Holger Drass } & {\small Bruno Quint } &
\begin{minipage}[]{6cm}
\smallskip
{\small Integration specialist\\
Systems Engineer }
\medskip
\end{minipage}
\\ \hline
\href{https://jira.lsstcorp.org/secure/Tests.jspa#/testCase/LVV-T2214}{LVV-T2214}
& {\small Holger Drass } & {\small  } &
\begin{minipage}[]{6cm}
\smallskip
{\small Systems Engineer\\
Integration specialist }
\medskip
\end{minipage}
\\ \hline
\href{https://jira.lsstcorp.org/secure/Tests.jspa#/testCase/LVV-T2215}{LVV-T2215}
& {\small Holger Drass } & {\small  } &
\begin{minipage}[]{6cm}
\smallskip
{\small Systems Verification Engineer\\
System Integration Specialist }
\medskip
\end{minipage}
\\ \hline
\href{https://jira.lsstcorp.org/secure/Tests.jspa#/testCase/LVV-T2216}{LVV-T2216}
& {\small Holger Drass } & {\small  } &
\begin{minipage}[]{6cm}
\smallskip
{\small (1) Integration Specialist\\
(1) Systems Engineer }
\medskip
\end{minipage}
\\ \hline
\href{https://jira.lsstcorp.org/secure/Tests.jspa#/testCase/LVV-T2241}{LVV-T2241}
& {\small Holger Drass } & {\small  } &
\begin{minipage}[]{6cm}
\smallskip
{\small Integration Specialist\\
Systems Engineer }
\medskip
\end{minipage}
\\ \hline
\href{https://jira.lsstcorp.org/secure/Tests.jspa#/testCase/LVV-T2242}{LVV-T2242}
& {\small Holger Drass } & {\small  } &
\begin{minipage}[]{6cm}
\smallskip
{\small Integration Specialist\\
Systems Engineer }
\medskip
\end{minipage}
\\ \hline
\multicolumn{2}{r}{T. Cycle \href{https://jira.lsstcorp.org/secure/Tests.jspa\#/testCycle/LVV-C174}{LVV-C174} owner:} &
\multicolumn{2}{l}{\textbf{
Holger Drass }
} \\\hline
\textbf{Test Cases} & \textbf{Assigned to} & \textbf{Executed by} & \textbf{Additional Test Personnel} \\ \hline
\href{https://jira.lsstcorp.org/secure/Tests.jspa#/testCase/LVV-T2193}{LVV-T2193}
& {\small Bo Xin [X] } & {\small  } &
\begin{minipage}[]{6cm}
\smallskip
{\small  }
\medskip
\end{minipage}
\\ \hline
\href{https://jira.lsstcorp.org/secure/Tests.jspa#/testCase/LVV-T2242}{LVV-T2242}
& {\small Bo Xin [X] } & {\small  } &
\begin{minipage}[]{6cm}
\smallskip
{\small Integration Specialist\\
Systems Engineer }
\medskip
\end{minipage}
\\ \hline
\href{https://jira.lsstcorp.org/secure/Tests.jspa#/testCase/LVV-T2213}{LVV-T2213}
& {\small Austin Roberts } & {\small Bruno Quint } &
\begin{minipage}[]{6cm}
\smallskip
{\small Integration specialist\\
Systems Engineer }
\medskip
\end{minipage}
\\ \hline
\href{https://jira.lsstcorp.org/secure/Tests.jspa#/testCase/LVV-T2190}{LVV-T2190}
& {\small Bo Xin [X] } & {\small  } &
\begin{minipage}[]{6cm}
\smallskip
{\small Systems Integration Specialist\\
Systems Verification Engineer }
\medskip
\end{minipage}
\\ \hline
\href{https://jira.lsstcorp.org/secure/Tests.jspa#/testCase/LVV-T2241}{LVV-T2241}
& {\small Bo Xin [X] } & {\small  } &
\begin{minipage}[]{6cm}
\smallskip
{\small Integration Specialist\\
Systems Engineer }
\medskip
\end{minipage}
\\ \hline
\href{https://jira.lsstcorp.org/secure/Tests.jspa#/testCase/LVV-T2214}{LVV-T2214}
& {\small Austin Roberts } & {\small  } &
\begin{minipage}[]{6cm}
\smallskip
{\small Systems Engineer\\
Integration specialist }
\medskip
\end{minipage}
\\ \hline
\href{https://jira.lsstcorp.org/secure/Tests.jspa#/testCase/LVV-T2215}{LVV-T2215}
& {\small Austin Roberts } & {\small  } &
\begin{minipage}[]{6cm}
\smallskip
{\small Systems Verification Engineer\\
System Integration Specialist }
\medskip
\end{minipage}
\\ \hline
\href{https://jira.lsstcorp.org/secure/Tests.jspa#/testCase/LVV-T2216}{LVV-T2216}
& {\small Austin Roberts } & {\small  } &
\begin{minipage}[]{6cm}
\smallskip
{\small (1) Integration Specialist\\
(1) Systems Engineer }
\medskip
\end{minipage}
\\ \hline
\multicolumn{2}{r}{T. Cycle \href{https://jira.lsstcorp.org/secure/Tests.jspa\#/testCycle/LVV-C175}{LVV-C175} owner:} &
\multicolumn{2}{l}{\textbf{
Holger Drass }
} \\\hline
\textbf{Test Cases} & \textbf{Assigned to} & \textbf{Executed by} & \textbf{Additional Test Personnel} \\ \hline
\href{https://jira.lsstcorp.org/secure/Tests.jspa#/testCase/LVV-T2215}{LVV-T2215}
& {\small Austin Roberts } & {\small Sandrine Thomas } &
\begin{minipage}[]{6cm}
\smallskip
{\small Systems Verification Engineer\\
System Integration Specialist }
\medskip
\end{minipage}
\\ \hline
\href{https://jira.lsstcorp.org/secure/Tests.jspa#/testCase/LVV-T2216}{LVV-T2216}
& {\small Austin Roberts } & {\small Sandrine Thomas } &
\begin{minipage}[]{6cm}
\smallskip
{\small (1) Integration Specialist\\
(1) Systems Engineer }
\medskip
\end{minipage}
\\ \hline
\multicolumn{2}{r}{T. Cycle \href{https://jira.lsstcorp.org/secure/Tests.jspa\#/testCycle/LVV-C176}{LVV-C176} owner:} &
\multicolumn{2}{l}{\textbf{
Holger Drass }
} \\\hline
\textbf{Test Cases} & \textbf{Assigned to} & \textbf{Executed by} & \textbf{Additional Test Personnel} \\ \hline
\href{https://jira.lsstcorp.org/secure/Tests.jspa#/testCase/LVV-T2578}{LVV-T2578}
& {\small Brian Stalder } & {\small  } &
\begin{minipage}[]{6cm}
\smallskip
{\small Systems Integration Specialist\\
Systems Verification Engineer }
\medskip
\end{minipage}
\\ \hline
\href{https://jira.lsstcorp.org/secure/Tests.jspa#/testCase/LVV-T2344}{LVV-T2344}
& {\small Bruno Quint } & {\small Holger Drass } &
\begin{minipage}[]{6cm}
\smallskip
{\small Systems Integration Specialist\\
Systems Verification Engineer }
\medskip
\end{minipage}
\\ \hline
\href{https://jira.lsstcorp.org/secure/Tests.jspa#/testCase/LVV-T2190}{LVV-T2190}
& {\small Bruno Quint } & {\small  } &
\begin{minipage}[]{6cm}
\smallskip
{\small Systems Integration Specialist\\
Systems Verification Engineer }
\medskip
\end{minipage}
\\ \hline
\href{https://jira.lsstcorp.org/secure/Tests.jspa#/testCase/LVV-T2290}{LVV-T2290}
& {\small Bruno Quint } & {\small Holger Drass } &
\begin{minipage}[]{6cm}
\smallskip
{\small Observing Specialist\\
Software Engineer\\
Systems Engineer }
\medskip
\end{minipage}
\\ \hline
\href{https://jira.lsstcorp.org/secure/Tests.jspa#/testCase/LVV-T2193}{LVV-T2193}
& {\small Bruno Quint } & {\small Holger Drass } &
\begin{minipage}[]{6cm}
\smallskip
{\small  }
\medskip
\end{minipage}
\\ \hline
\href{https://jira.lsstcorp.org/secure/Tests.jspa#/testCase/LVV-T2228}{LVV-T2228}
& {\small Bruno Quint } & {\small Bruno Quint } &
\begin{minipage}[]{6cm}
\smallskip
{\small Systems Integration Specialist\\
Systems Verification Engineer }
\medskip
\end{minipage}
\\ \hline
\href{https://jira.lsstcorp.org/secure/Tests.jspa#/testCase/LVV-T2229}{LVV-T2229}
& {\small Bruno Quint } & {\small  } &
\begin{minipage}[]{6cm}
\smallskip
{\small Systems Integration Specialist\\
Systems Verification Engineer\\
Software Engineer }
\medskip
\end{minipage}
\\ \hline
\multicolumn{2}{r}{T. Cycle \href{https://jira.lsstcorp.org/secure/Tests.jspa\#/testCycle/LVV-C177}{LVV-C177} owner:} &
\multicolumn{2}{l}{\textbf{
Holger Drass }
} \\\hline
\textbf{Test Cases} & \textbf{Assigned to} & \textbf{Executed by} & \textbf{Additional Test Personnel} \\ \hline
\href{https://jira.lsstcorp.org/secure/Tests.jspa#/testCase/LVV-T2540}{LVV-T2540}
& {\small Austin Roberts } & {\small  } &
\begin{minipage}[]{6cm}
\smallskip
{\small (1) Systems Engineer, (1) Electrical Engineer }
\medskip
\end{minipage}
\\ \hline
\href{https://jira.lsstcorp.org/secure/Tests.jspa#/testCase/LVV-T2227}{LVV-T2227}
& {\small Austin Roberts } & {\small  } &
\begin{minipage}[]{6cm}
\smallskip
{\small (1) Systems Engineer, (1) Electrical Engineer }
\medskip
\end{minipage}
\\ \hline
\href{https://jira.lsstcorp.org/secure/Tests.jspa#/testCase/LVV-T2231}{LVV-T2231}
& {\small Austin Roberts } & {\small  } &
\begin{minipage}[]{6cm}
\smallskip
{\small (1) Systems Engineer, (1) Electrical Engineer }
\medskip
\end{minipage}
\\ \hline
\multicolumn{2}{r}{T. Cycle \href{https://jira.lsstcorp.org/secure/Tests.jspa\#/testCycle/LVV-C206}{LVV-C206} owner:} &
\multicolumn{2}{l}{\textbf{
Holger Drass }
} \\\hline
\textbf{Test Cases} & \textbf{Assigned to} & \textbf{Executed by} & \textbf{Additional Test Personnel} \\ \hline
\href{https://jira.lsstcorp.org/secure/Tests.jspa#/testCase/LVV-T2215}{LVV-T2215}
& {\small Sandrine Thomas } & {\small  } &
\begin{minipage}[]{6cm}
\smallskip
{\small Systems Verification Engineer\\
System Integration Specialist }
\medskip
\end{minipage}
\\ \hline
\href{https://jira.lsstcorp.org/secure/Tests.jspa#/testCase/LVV-T2216}{LVV-T2216}
& {\small Sandrine Thomas } & {\small  } &
\begin{minipage}[]{6cm}
\smallskip
{\small (1) Integration Specialist\\
(1) Systems Engineer }
\medskip
\end{minipage}
\\ \hline
\href{https://jira.lsstcorp.org/secure/Tests.jspa#/testCase/LVV-T2241}{LVV-T2241}
& {\small Bo Xin [X] } & {\small Chris Walter } &
\begin{minipage}[]{6cm}
\smallskip
{\small Integration Specialist\\
Systems Engineer }
\medskip
\end{minipage}
\\ \hline
\href{https://jira.lsstcorp.org/secure/Tests.jspa#/testCase/LVV-T2213}{LVV-T2213}
& {\small Austin Roberts } & {\small Ioana Sotuela } &
\begin{minipage}[]{6cm}
\smallskip
{\small Integration specialist\\
Systems Engineer }
\medskip
\end{minipage}
\\ \hline
\href{https://jira.lsstcorp.org/secure/Tests.jspa#/testCase/LVV-T2214}{LVV-T2214}
& {\small Austin Roberts } & {\small Ioana Sotuela } &
\begin{minipage}[]{6cm}
\smallskip
{\small Systems Engineer\\
Integration specialist }
\medskip
\end{minipage}
\\ \hline
\end{longtable}
}

\newpage

\section{Test Campaign Overview}
\label{sect:overview}

\subsection{Summary}
\label{sect:summarytable}

{\small
\begin{longtable}{p{2cm}cp{2.3cm}p{8.6cm}p{2.3cm}}
\toprule
\multicolumn{2}{r}{ T. Plan \href{https://jira.lsstcorp.org/secure/Tests.jspa\#/testPlan/LVV-P81}{LVV-P81}:} &
\multicolumn{2}{p{10.9cm}}{\textbf{ Level 3 System Spread Configuration Integration }} & Approved \\\hline
\multicolumn{2}{r}{ T. Cycle \href{https://jira.lsstcorp.org/secure/Tests.jspa\#/testCycle/LVV-C170}{LVV-C170}:} &
\multicolumn{2}{p{10.9cm}}{\textbf{ Level 3 System Integration: Individual Subsystem Interface to TCS }} & In Progress \\\hline
\textbf{Test Cases} &  \textbf{Ver.}  \\\toprule
\href{https://jira.lsstcorp.org/secure/Tests.jspa#/testCase/LVV-T1802}{LVV-T1802}
&  2
\\
\href{https://jira.lsstcorp.org/secure/Tests.jspa#/testCase/LVV-T1615}{LVV-T1615}
&  2
\\
\href{https://jira.lsstcorp.org/secure/Tests.jspa#/testCase/LVV-T1600}{LVV-T1600}
&  2
\\
\href{https://jira.lsstcorp.org/secure/Tests.jspa#/testCase/LVV-T1602}{LVV-T1602}
&  3
\\
\href{https://jira.lsstcorp.org/secure/Tests.jspa#/testCase/LVV-T2222}{LVV-T2222}
&  1
\\
\href{https://jira.lsstcorp.org/secure/Tests.jspa#/testCase/LVV-T2221}{LVV-T2221}
&  2
\\
\href{https://jira.lsstcorp.org/secure/Tests.jspa#/testCase/LVV-T2232}{LVV-T2232}
&  2
\\
\href{https://jira.lsstcorp.org/secure/Tests.jspa#/testCase/LVV-T2219}{LVV-T2219}
&  1
\\
\\\hline
\multicolumn{2}{r}{ T. Cycle \href{https://jira.lsstcorp.org/secure/Tests.jspa\#/testCycle/LVV-C172}{LVV-C172}:} &
\multicolumn{2}{p{10.9cm}}{\textbf{ Level 3 System Integration: M2 Hardware with M1M3 Hardware and
Simulators }} & Not Executed \\\hline
\textbf{Test Cases} &  \textbf{Ver.}  \\\toprule
\href{https://jira.lsstcorp.org/secure/Tests.jspa#/testCase/LVV-T2193}{LVV-T2193}
&  1
\\
\href{https://jira.lsstcorp.org/secure/Tests.jspa#/testCase/LVV-T2242}{LVV-T2242}
&  1
\\
\href{https://jira.lsstcorp.org/secure/Tests.jspa#/testCase/LVV-T2213}{LVV-T2213}
&  1
\\
\href{https://jira.lsstcorp.org/secure/Tests.jspa#/testCase/LVV-T2190}{LVV-T2190}
&  1
\\
\href{https://jira.lsstcorp.org/secure/Tests.jspa#/testCase/LVV-T2241}{LVV-T2241}
&  1
\\
\href{https://jira.lsstcorp.org/secure/Tests.jspa#/testCase/LVV-T2214}{LVV-T2214}
&  1
\\
\href{https://jira.lsstcorp.org/secure/Tests.jspa#/testCase/LVV-T2215}{LVV-T2215}
&  2
\\
\href{https://jira.lsstcorp.org/secure/Tests.jspa#/testCase/LVV-T2216}{LVV-T2216}
&  1
\\
\\\hline
\multicolumn{2}{r}{ T. Cycle \href{https://jira.lsstcorp.org/secure/Tests.jspa\#/testCycle/LVV-C173}{LVV-C173}:} &
\multicolumn{2}{p{10.9cm}}{\textbf{ Level 3 System Integration: M1M3 Hardware with Simulators }} & In Progress \\\hline
\textbf{Test Cases} &  \textbf{Ver.}  \\\toprule
\href{https://jira.lsstcorp.org/secure/Tests.jspa#/testCase/LVV-T2190}{LVV-T2190}
&  1
\\
\href{https://jira.lsstcorp.org/secure/Tests.jspa#/testCase/LVV-T2193}{LVV-T2193}
&  1
\\
\href{https://jira.lsstcorp.org/secure/Tests.jspa#/testCase/LVV-T2213}{LVV-T2213}
&  1
\\
\href{https://jira.lsstcorp.org/secure/Tests.jspa#/testCase/LVV-T2214}{LVV-T2214}
&  1
\\
\href{https://jira.lsstcorp.org/secure/Tests.jspa#/testCase/LVV-T2215}{LVV-T2215}
&  2
\\
\href{https://jira.lsstcorp.org/secure/Tests.jspa#/testCase/LVV-T2216}{LVV-T2216}
&  1
\\
\href{https://jira.lsstcorp.org/secure/Tests.jspa#/testCase/LVV-T2241}{LVV-T2241}
&  1
\\
\href{https://jira.lsstcorp.org/secure/Tests.jspa#/testCase/LVV-T2242}{LVV-T2242}
&  1
\\
\\\hline
\multicolumn{2}{r}{ T. Cycle \href{https://jira.lsstcorp.org/secure/Tests.jspa\#/testCycle/LVV-C174}{LVV-C174}:} &
\multicolumn{2}{p{10.9cm}}{\textbf{ Level 3 System Integration: Camera Hexapod Hardware with M2 Hardware,
M1M3 Hardware, and Simulators }} & In Progress \\\hline
\textbf{Test Cases} &  \textbf{Ver.}  \\\toprule
\href{https://jira.lsstcorp.org/secure/Tests.jspa#/testCase/LVV-T2193}{LVV-T2193}
&  1
\\
\href{https://jira.lsstcorp.org/secure/Tests.jspa#/testCase/LVV-T2242}{LVV-T2242}
&  1
\\
\href{https://jira.lsstcorp.org/secure/Tests.jspa#/testCase/LVV-T2213}{LVV-T2213}
&  1
\\
\href{https://jira.lsstcorp.org/secure/Tests.jspa#/testCase/LVV-T2190}{LVV-T2190}
&  1
\\
\href{https://jira.lsstcorp.org/secure/Tests.jspa#/testCase/LVV-T2241}{LVV-T2241}
&  1
\\
\href{https://jira.lsstcorp.org/secure/Tests.jspa#/testCase/LVV-T2214}{LVV-T2214}
&  1
\\
\href{https://jira.lsstcorp.org/secure/Tests.jspa#/testCase/LVV-T2215}{LVV-T2215}
&  2
\\
\href{https://jira.lsstcorp.org/secure/Tests.jspa#/testCase/LVV-T2216}{LVV-T2216}
&  1
\\
\\\hline
\multicolumn{2}{r}{ T. Cycle \href{https://jira.lsstcorp.org/secure/Tests.jspa\#/testCycle/LVV-C175}{LVV-C175}:} &
\multicolumn{2}{p{10.9cm}}{\textbf{ Level 3 System Integration: Camera Rotator Hardware with Camera Hexapod
Hardware, M2 Hardware, M1M3 Hardware, and Simulators }} & In Progress \\\hline
\textbf{Test Cases} &  \textbf{Ver.}  \\\toprule
\href{https://jira.lsstcorp.org/secure/Tests.jspa#/testCase/LVV-T2215}{LVV-T2215}
&  2
\\
\href{https://jira.lsstcorp.org/secure/Tests.jspa#/testCase/LVV-T2216}{LVV-T2216}
&  1
\\
\\\hline
\multicolumn{2}{r}{ T. Cycle \href{https://jira.lsstcorp.org/secure/Tests.jspa\#/testCycle/LVV-C176}{LVV-C176}:} &
\multicolumn{2}{p{10.9cm}}{\textbf{ Level 3 System Integration: ComCam Software with Camera Rotator
Hardware, Camera Hexapod Hardware, M2 Hardware, M1M3 Hardware, and
Simulators }} & In Progress \\\hline
\textbf{Test Cases} &  \textbf{Ver.}  \\\toprule
\href{https://jira.lsstcorp.org/secure/Tests.jspa#/testCase/LVV-T2578}{LVV-T2578}
&  1
\\
\href{https://jira.lsstcorp.org/secure/Tests.jspa#/testCase/LVV-T2344}{LVV-T2344}
&  1
\\
\href{https://jira.lsstcorp.org/secure/Tests.jspa#/testCase/LVV-T2190}{LVV-T2190}
&  1
\\
\href{https://jira.lsstcorp.org/secure/Tests.jspa#/testCase/LVV-T2290}{LVV-T2290}
&  2
\\
\href{https://jira.lsstcorp.org/secure/Tests.jspa#/testCase/LVV-T2193}{LVV-T2193}
&  1
\\
\href{https://jira.lsstcorp.org/secure/Tests.jspa#/testCase/LVV-T2228}{LVV-T2228}
&  1
\\
\href{https://jira.lsstcorp.org/secure/Tests.jspa#/testCase/LVV-T2229}{LVV-T2229}
&  2
\\
\\\hline
\multicolumn{2}{r}{ T. Cycle \href{https://jira.lsstcorp.org/secure/Tests.jspa\#/testCycle/LVV-C177}{LVV-C177}:} &
\multicolumn{2}{p{10.9cm}}{\textbf{ Level 3 System Integration: GIS Integration }} & Not Executed \\\hline
\textbf{Test Cases} &  \textbf{Ver.}  \\\toprule
\href{https://jira.lsstcorp.org/secure/Tests.jspa#/testCase/LVV-T2540}{LVV-T2540}
&  1
\\
\href{https://jira.lsstcorp.org/secure/Tests.jspa#/testCase/LVV-T2227}{LVV-T2227}
&  1
\\
\href{https://jira.lsstcorp.org/secure/Tests.jspa#/testCase/LVV-T2231}{LVV-T2231}
&  1
\\
\\\hline
\multicolumn{2}{r}{ T. Cycle \href{https://jira.lsstcorp.org/secure/Tests.jspa\#/testCycle/LVV-C206}{LVV-C206}:} &
\multicolumn{2}{p{10.9cm}}{\textbf{ Level 3 System Integration: M2 Hardware, M1M3 Hardware, CCW+Rotator
Hardware, Cam Hex and Simulators }} & In Progress \\\hline
\textbf{Test Cases} &  \textbf{Ver.}  \\\toprule
\href{https://jira.lsstcorp.org/secure/Tests.jspa#/testCase/LVV-T2215}{LVV-T2215}
&  2
\\
\href{https://jira.lsstcorp.org/secure/Tests.jspa#/testCase/LVV-T2216}{LVV-T2216}
&  1
\\
\href{https://jira.lsstcorp.org/secure/Tests.jspa#/testCase/LVV-T2241}{LVV-T2241}
&  1
\\
\href{https://jira.lsstcorp.org/secure/Tests.jspa#/testCase/LVV-T2213}{LVV-T2213}
&  1
\\
\href{https://jira.lsstcorp.org/secure/Tests.jspa#/testCase/LVV-T2214}{LVV-T2214}
&  1
\\
\\\hline
\caption{Test Campaign Summary}
\label{table:summary}
\end{longtable}
}

\subsection{Overall Assessment}
\label{sect:overallassessment}

Not yet available.

\subsection{Recommended Improvements}
\label{sect:recommendations}

\newpage
\section{Detailed Tests}
\label{sect:detailedtests}

\subsection{Test Cycle LVV-C170 }

Open test cycle {\it \href{https://jira.lsstcorp.org/secure/Tests.jspa#/testrun/LVV-C170}{Level 3 System Integration: Individual Subsystem Interface to TCS}} in Jira.

Test Cycle name: Level 3 System Integration: Individual Subsystem Interface to TCS\\
Status: In Progress

Verify the interface with SAL for each of the following subsystems:

\begin{itemize}
\tightlist
\item
  Camera Hexapod per \citeds{LTS-160}
\item
  Camera Rotator per \citeds{LTS-160}
\item
  M2 per \citeds{LTS-162}
\item
  M1M3 per \citeds{LTS-161}
\item
  ComCam per XML
\item
  M2 Hexapod (simulator) per \citeds{LTS-160}
\item
  MTMount (simulator) per \citeds{LTS-159}
\item
  MTAOS per \citeds{LTS-163}
\end{itemize}

\subsubsection{Software Version/Baseline}
\begin{itemize}
\tightlist
\item
  XML 12.0 and Cycle 26 rev 0.
\item
  SAL object versions and CSC versions:
  https://github.com/lsst-ts/ts\_cycle\_build/blob/c0025.008/cycle/cycle.env
\item
  EUI versions can be determined from the Puppet.
\item
  At a later point, the controller version can be determined from
  Puppet. At the moment the controller versions can be determined from
  the Tags on GitHub from the deployed system itself.
\item
  MTHexapod\_1 (CamHex) low-level controller - v1.4.1
\item
  MTHexapod\_1 (CamHex) EUI - v1.3.2
\item
  {MTRot (Camera Rotator) EUI-~v1.2.2}
\item
  MTRot (Camera Rotator) low-level controller - v1.5.2
\item
  MTRot (Camera Rotator) Trajectory generator (Tekniker version)- v1.2.9
\item
  MTM2 low-level controller v1.0.0
\item
  MTM2 EUI -{~v1.4.0}
\item
  MTM1M3 low-level controller - {vX.X}
\item
  {MTM1M3 EUI -~}{vX.X}
\item
  ComCam low-level controller- {vX.X}
\item
  {ComCam CCS}{- vX.X}
\item
  {MTHexapod\_2 (M2 Hexapod) simulator~}{-vX.X}{\\
  }
\item
  MTMount simulator/hardware hybrid (CCW is real hardware) version -
  {vX.X}
\end{itemize}

\subsubsection{Configuration}
Each hardware subsystem and the simulators are tested independently from
each other. All hardware components are located on the level 3 floor.

\subsubsection{Test Cases in LVV-C170 Test Cycle}

\paragraph{ LVV-T1802 - Integration of M2 Hexapod with SAL }\mbox{}\\

Version \textbf{2}.
Open  \href{https://jira.lsstcorp.org/secure/Tests.jspa#/testCase/LVV-T1802}{\textit{ LVV-T1802 } }
test case in Jira.

The objective of this test case is to re-verify the functional
requirements of the M2 hexapod's software, after shipment of the
hardware from the vendor's facility to the Summit, as defined in \citeds{LTS-206}
and \citeds{LTS-160}. This test case will only exercise the functionality that
was executed previously and meets the following criteria:

\begin{itemize}
\tightlist
\item
  Only requires the use of Rubin Observatory code to replace MOOG's
  middleware code
\item
  Only requires the M2 hexapod to be operable
\item
  Only requires command through the CSC after the PXI real-time
  controller is switched from GUI mode to DDS mode
\item
  Only requires testing the synchronous mode

  \begin{itemize}
  \tightlist
  \item
    \textbf{Asynchronous mode is not a standard mode of operation}
  \end{itemize}
\item
  Does require the M2 hexapod temperature sensors to be operating
\item
  Does \textbf{NOT} require the M2 hexapod to be rotated to various
  elevation angles
\item
  Does \textbf{NOT~}require the M2 hexapod to be in a climate-controlled
  environment
\end{itemize}

The software functional requirements were previously verified during the
test campaign by the vendor at the vendor's facility and accepted by
Rubin Observatory during the Factory Acceptance Test review. The test
procedure used during the vendor's acceptance testing is the \emph{LSST
Hexapods-Rotator Software Acceptance Test Procedure} which is attached
to this test case. The test steps of this test case are the same steps
from the procedure for the testing of the Camera Hexapod. The order of
the steps was changed to reflect the \emph{Proposal of Hexapod Test~on
the Dec. 2019~}Confluence page which can be found linked in the
Traceability tab.\\[2\baselineskip]See the attached \emph{LSST Hexapod
Operator's Manual} for more information on how to operate the hexapod.

\textbf{ Preconditions}:\\
Prior to the execution of this test case to re-verify the M2 Hexapod
hardware functional requirements, the following Summit tasks must be
completed:

\begin{itemize}
\tightlist
\item
  The measurement equipment has been set-up for testing

  \begin{itemize}
  \tightlist
  \item
    \url{https://jira.lsstcorp.org/browse/SUMMIT-1943}
  \end{itemize}
\end{itemize}

\paragraph{ LVV-T1615 - Integration of M2 with SAL }\mbox{}\\

Version \textbf{2}.
Open  \href{https://jira.lsstcorp.org/secure/Tests.jspa#/testCase/LVV-T1615}{\textit{ LVV-T1615 } }
test case in Jira.

The objective of this test case is to verify the software requirements
of the M2 cell, as defined in \citeds{LTS-146} and \citeds{LTS-162}, utilizing the updated
M2 CSC (SAL). The software requirements were previously verified during
the test campaign by the vendor, but have since been updated per
CR-0202. This test case will exercise the functionality of the M2 CSC
and meets the following criteria:

\begin{itemize}
\tightlist
\item
  Only requires the most current version of SAL
\item
  Only requires the M2 surrogate to be loaded on the cell
\item
  Only requires the use of the DDS
\item
  Does \textbf{NOT} require the M2 to be integrated with the Interlock
  system
\end{itemize}

\textbf{ Preconditions}:\\


\paragraph{ LVV-T1600 - Integration of Camera Hexapod with SAL }\mbox{}\\

Version \textbf{2}.
Open  \href{https://jira.lsstcorp.org/secure/Tests.jspa#/testCase/LVV-T1600}{\textit{ LVV-T1600 } }
test case in Jira.

The objective of this test case is to re-verify the functional
requirements of the camera hexapod's software, after shipment of the
hardware from the vendor's facility to the Summit, as defined in \citeds{LTS-206}
and \citeds{LTS-160}. This test case will only exercise the functionality that
was executed previously and meets the following criteria:

\begin{itemize}
\tightlist
\item
  Only requires the use of Russell's code to replace MOOG's middleware
  code
\item
  Only requires the camera hexapod to be operable
\item
  Only requires command through the CSC after the cRIO is switched from
  GUI mode to DDS mode
\item
  Only requires testing of the synchronous mode

  \begin{itemize}
  \tightlist
  \item
    \textbf{Asynchronous mode is not a standard mode of operation}
  \end{itemize}
\item
  This test case can be executed with or without the camera rotator to
  be loaded with the camera simulated mass or actual camera hardware
\end{itemize}

The software functional requirements were previously verified during the
test campaign by the vendor at the vendor's facility and accepted by
LSST during the Factory Acceptance Test review. The test procedure used
during the vendor's acceptance testing is the \emph{LSST
Hexapods-Rotator Software Acceptance Test Procedure} which is attached
to this test case. The test steps of this test case are derived from the
same procedure, but the order of the steps have been changed to reflect
the \emph{Proposal of Hexapod Test~on Dec. 2019~}Confluence page which
can be found linked in the Traceability tab.\\[2\baselineskip]See the
attached \emph{LSST Rotator Hexapod's Manual} for more information on
how to operate the hexapod.

\textbf{ Preconditions}:\\
Prior to the execution of this test case to re-verify the Camera Hexapod
hardware functional requirements, the following Summit tasks must be
completed:

\begin{itemize}
\tightlist
\item
  The Hexapod has been installed on the camera cart

  \begin{itemize}
  \tightlist
  \item
    \url{https://jira.lsstcorp.org/browse/SUMMIT-3224}
  \end{itemize}
\item
  The Hexapod Controller has been deployed on the summit

  \begin{itemize}
  \tightlist
  \item
    \url{https://jira.lsstcorp.org/browse/SUMMIT-3229}
  \end{itemize}
\item
  Boxes for the Hexapod have been transported to the 3rd level

  \begin{itemize}
  \tightlist
  \item
    \url{https://jira.lsstcorp.org/browse/SUMMIT-3230}
  \end{itemize}
\item
  All Hexapod cables and cabinets have been prepared for integration
  with camera cart

  \begin{itemize}
  \tightlist
  \item
    \url{https://jira.lsstcorp.org/browse/SUMMIT-3231}
  \end{itemize}
\item
  The offset has been installed onto the integrating structure

  \begin{itemize}
  \tightlist
  \item
    \url{https://jira.lsstcorp.org/browse/SUMMIT-3293}
  \end{itemize}
\item
  The Camera Hexapod electrical connections have been tested

  \begin{itemize}
  \tightlist
  \item
    \url{https://jira.lsstcorp.org/browse/SUMMIT-3294}
  \end{itemize}
\end{itemize}

\paragraph{ LVV-T1602 - Integration of Camera Rotator with SAL }\mbox{}\\

Version \textbf{3}.
Open  \href{https://jira.lsstcorp.org/secure/Tests.jspa#/testCase/LVV-T1602}{\textit{ LVV-T1602 } }
test case in Jira.

The objective of this test case is to verify the software requirements
of the camera rotator, as defined in \citeds{LTS-160}.\\
This test case will only exercise the functionality that was executed
previously and meets the following criteria:

\begin{itemize}
\tightlist
\item
  Only requires the camera rotator to be operable
\item
  Only requires command through the CSC after the cRIO is switched from
  GUI mode to DDS mode
\item
  Require the rotator to be loaded with the camera simulated mass or
  actual camera hardware
\end{itemize}

The test procedure used during the previous testing without SAL was the
\emph{LSST Hexapods-Rotator Software Acceptance Test Procedure.~}\\
The test steps of this test case are taken directly from that document
on how to perform the test in a similar way with the use of
SAL.\\[2\baselineskip]History of this test case:\\
See the \emph{New Rotator CSC Test at SLAC on 10/17/2019~}linked in the
Traceability tab for the results of Russell's test of the events and
telemetry with the bare rotator in SLAC.\\[2\baselineskip]See the
attached \emph{LSST Rotator Operator's Manual} for more information on
how to operate the rotator.\\[2\baselineskip]This Version 3.0 is
detailed for XML 12.0 and Cycle 26.\\[2\baselineskip]While executing
this test case ensure that the values for events and telemetry are
reasonable.

\textbf{ Preconditions}:\\
Prior to the execution of this test case, the following Summit tasks
must be completed:\\[2\baselineskip]

\begin{itemize}
\tightlist
\item
  The ComCam has been installed on the Camera Rotator

  \begin{itemize}
  \tightlist
  \item
    \url{https://jira.lsstcorp.org/browse/SUMMIT-4935}
  \end{itemize}
\item
  The ComCam utility lines have been connected to the CCW~

  \begin{itemize}
  \tightlist
  \item
    h\url{https://jira.lsstcorp.org/browse/SUMMIT-4936}
  \end{itemize}
\item
  The functional testing of the Camera Rotator Hardware~

  \begin{itemize}
  \tightlist
  \item
    h\href{https://jira.lsstcorp.org/browse/SUMMIT-3370}{ttps://jira.lsstcorp.org/browse/SUMMIT-3370}
  \end{itemize}
\item
  The functional testing of the Camera Rotator Software without SAL

  \begin{itemize}
  \tightlist
  \item
    \url{https://jira.lsstcorp.org/browse/SUMMIT-3371}
  \end{itemize}
\end{itemize}

\paragraph{ LVV-T2222 - MTAOS Interface to SAL }\mbox{}\\

Version \textbf{1}.
Open  \href{https://jira.lsstcorp.org/secure/Tests.jspa#/testCase/LVV-T2222}{\textit{ LVV-T2222 } }
test case in Jira.

The objective of this test case is to verify the MTAOS communication
interface with SAL by checking the commands as defined in the latest
XML.\\
This test will not verify the functionality of the commands, but will
simply test that the MTAOS is able to publish the commands to SAL.\\
The MTAOS commands, as well as the events, will be fully verified in
later test cycles of the Level 3 Integration Test Plan when the MTAOS is
integrated with the M1M3, M2, and the hexapods.

\textbf{ Preconditions}:\\
\begin{itemize}
\tightlist
\item
  Only requires the MTAOS to be operable~
\item
  Requires the use of SAL/DDS and the EFD
\item
  Does not require any communication with the M1M3, M2, and the hexapods
\end{itemize}

\paragraph{ LVV-T2221 - MTMount Interface to SAL }\mbox{}\\

Version \textbf{2}.
Open  \href{https://jira.lsstcorp.org/secure/Tests.jspa#/testCase/LVV-T2221}{\textit{ LVV-T2221 } }
test case in Jira.

The objective of this test case is to verify~

\begin{itemize}
\tightlist
\item
  the MTMount communication interface with SAL by testing the commands,
  events and telemetry defined by the latest version of the XML
\item
  all topic parameters are filled with data
\item
  the data are reasonable.
\end{itemize}

This test case will exercise the functionality of the Telescope Mount
and meet the following criteria:\\

\begin{itemize}
\tightlist
\item
  Only requires the most current version of SAL
\item
  Requires the use of the DDS and the EFD
\end{itemize}

\textbf{ Preconditions}:\\
No specific preconditions identified

\paragraph{ LVV-T2232 - M1M3 Integration with SAL }\mbox{}\\

Version \textbf{2}.
Open  \href{https://jira.lsstcorp.org/secure/Tests.jspa#/testCase/LVV-T2232}{\textit{ LVV-T2232 } }
test case in Jira.

The objective of this test case is to verify the latest M1M3 commands,
events and telemetry defined by the latest version of the XML. This test
case will exercise the functionality of the M1M3 on the 3rd level of the
Summit and meets the following criteria:

\begin{itemize}
\tightlist
\item
  Only requires the most current version of SAL
\item
  Only requires the M1M3 surrogate to be loaded on the cell
\item
  Requires the use of the DDS and the EFD
\end{itemize}

\textbf{ Preconditions}:\\
M1M3 minimal functionality testing passed successfully.

\paragraph{ LVV-T2219 - ComCam Interface to SAL }\mbox{}\\

Version \textbf{1}.
Open  \href{https://jira.lsstcorp.org/secure/Tests.jspa#/testCase/LVV-T2219}{\textit{ LVV-T2219 } }
test case in Jira.

The objective of this test case is to verify the ComCam communication
interface with SAL by testing the commands, events and telemetry defined
by the latest version of the XML. This test case will exercise the
functionality of the ComCam on the 3rd level of the Summit and meets the
following criteria:

\begin{itemize}
\tightlist
\item
  Only requires the ComCam to be operational
\item
  Does not require any communication with other subsystems
\item
  Does require the use of SAL and the EFD
\end{itemize}

\textbf{ Preconditions}:\\


\subsection{Test Cycle LVV-C172 }

Open test cycle {\it \href{https://jira.lsstcorp.org/secure/Tests.jspa#/testrun/LVV-C172}{Level 3 System Integration: M2 Hardware with M1M3 Hardware and
Simulators}} in Jira.

Test Cycle name: Level 3 System Integration: M2 Hardware with M1M3 Hardware and
Simulators\\
Status: Not Executed

Get details from Google doc.\\[3\baselineskip]This uses the MTMount
simulator.

\subsubsection{Software Version/Baseline}
Not provided.

\subsubsection{Configuration}
\begin{itemize}
\tightlist
\item
  GIS hardware
\item
  M1M3 hardware
\item
  M2 hardware
\item
  Camera Hexapod simulator
\item
  Camera Rotator simulator
\item
  M2 Hexapod simulator
\item
  MTMount simulator
\item
  MTAOS software
\item
  ComCam hardware but not utilizing software
\end{itemize}

\subsubsection{Test Cases in LVV-C172 Test Cycle}

\paragraph{ LVV-T2193 - MTAOS handling of rejected commands }\mbox{}\\

Version \textbf{1}.
Open  \href{https://jira.lsstcorp.org/secure/Tests.jspa#/testCase/LVV-T2193}{\textit{ LVV-T2193 } }
test case in Jira.

Testing the communications between MTAOS and M1M3, M2, and the hexapods,
when one set of MTAOS corrections gets rejected.\\
This test exercises the following components:

\begin{itemize}
\tightlist
\item
  M1M3
\item
  M2
\item
  Camera Hexapod
\item
  M2 Hexapod
\end{itemize}

Extention of LVV-T2190. It has the same validation ticket traced to it.

\textbf{ Preconditions}:\\
\begin{itemize}
\tightlist
\item
  M1M3, M2, and Camera Hexapod pass the minimum functionality tests (we
  use M2 hexapod simulator as long as necessary)
\item
  Summit network and EFD are ready to record Camera Hexapods commands,
  events, and telemetry (use of domain 0)
\item
  Jupyter notebook for the test has been written and tested in the NCSA
  simulated environment
\item
  MTAOS CSC is set up and ready on the summit network
\end{itemize}

\paragraph{ LVV-T2242 - MTAOS configuration changes }\mbox{}\\

Version \textbf{1}.
Open  \href{https://jira.lsstcorp.org/secure/Tests.jspa#/testCase/LVV-T2242}{\textit{ LVV-T2242 } }
test case in Jira.

Testing the communications between MTAOS (both WEP and OFC) and M1M3,
M2, and the hexapods, via the MTAOS \emph{addAberration} command.\\
This test exercises the following components:

\begin{itemize}
\tightlist
\item
  M1M3
\item
  M2
\item
  camera hexapod
\item
  M2 hexapod
\item
  MTAOS
\end{itemize}

From
\href{https://docushare.lsst.org/docushare/dsweb/Get/LTS-186}{LTS-186 -
MTAOS: WEP and OFC Requirement Document}:\\
- Kallman filtering parameters\\
- Number of visits\\
- Detector(s) to be used*\\
- Location of the detector*\\
- Source selection parameters ({TBC})*\\
- Basis Set\\
-
\href{https://github.com/lsst-ts/ts_config_mttcs/blob/develop/MTAOS/v2/ofc/comcam/senM_9_19_50.yaml}{Sensitivity
Matrix}*\\
- Reconstructor Matrix\\
- Penalty Matrix\\
- Quality Matrix\\
- Gain\\[2\baselineskip]* These are the configurations that we can
verify on System Spread Integration Tests on Level 3.

\textbf{ Preconditions}:\\
\begin{itemize}
\tightlist
\item
  M1M3, M2, and Camera Hexapod pass the minimum functionality tests (we
  use the M2 hexapod simulator)
\item
  Summit network and EFD are ready to record CamHex commands, events,
  and telemetry (use of domain 0)
\item
  Jupyter notebook for the test has been written and tested in the test
  stand
\item
  MTAOS CSC is set up and ready on the summit network
\end{itemize}

\paragraph{ LVV-T2213 - Look-up Table Application from MTMount Elevation and Azimuth Changes }\mbox{}\\

Version \textbf{1}.
Open  \href{https://jira.lsstcorp.org/secure/Tests.jspa#/testCase/LVV-T2213}{\textit{ LVV-T2213 } }
test case in Jira.

Demonstrate that M1M3, M2, and the hexapods can subscribe to MTMount
elevation changes, and the LUT is applied according to the MTMount
elevation angle.\\
Testing that the components are subscribing to the mount
telemetry\\[2\baselineskip]\citeds{LTS-88} for M1M3\\
\citeds{LTS-146} for M2\\
\citeds{LTS-206} for the hexapods\\
XML for the Mount\\
\citeds{LTS-160}\\
\citeds{LTS-161}\\
\citeds{LTS-162}

\textbf{ Preconditions}:\\
The following components are ready for testing (depending on which test
cycle this is being executed in, each component is either a hardware
component or a simulator):\\

\begin{itemize}
\tightlist
\item
  M1M3
\item
  M2~
\item
  M2 Hexapod
\item
  Camera Hexapod
\item
  MTMount
\end{itemize}

\paragraph{ LVV-T2190 - MTAOS add aberrations to M1M3+M2+Hexapods }\mbox{}\\

Version \textbf{1}.
Open  \href{https://jira.lsstcorp.org/secure/Tests.jspa#/testCase/LVV-T2190}{\textit{ LVV-T2190 } }
test case in Jira.

Testing the communications between MTAOS and M1M3, M2, and the hexapods,
via the MTAOS \emph{addAberration~}command.\\
This test exercises the following components:

\begin{itemize}
\tightlist
\item
  M1M3
\item
  M2
\item
  Camera Hexapod
\item
  M2 Hexapod
\item
  MTAOS
\end{itemize}

These test steps are to test adding aberrations in the way of 1um*2um =
2um.\\
Same requirements as in
\href{https://jira.lsstcorp.org/secure/Tests.jspa\#/testCase/LVV-T2241}{LVV-2241}.
Different ways to verify for them.\\[2\baselineskip]Requirements are
from \citeds{LTS-186} (MTAOS), \citeds{LTS-206}, \citeds{LTS-88}, and \citeds{LTS-146} (MTAOS calculates and
publishes the offsets.)\\
MTAOS commands the offsets to the other components LTS- 160,161,162
(Components subscribe to the MTAOS telemetry and apply the offsets.)

\textbf{ Preconditions}:\\
\begin{itemize}
\tightlist
\item
  M1M3, M2, and Camera Hexapod pass the minimum functionality test (we
  use the M2 hexapod simulator)
\item
  Summit network and EFD are ready to record Camera Hexapod's commands,
  events, and telemetry (use of domain 0)
\item
  Jupyter notebook for the test has been written and tested a test stand
\item
  MTAOS CSC is set up and ready on the summit network
\end{itemize}

\paragraph{ LVV-T2241 - MTAOS corrections accumulation }\mbox{}\\

Version \textbf{1}.
Open  \href{https://jira.lsstcorp.org/secure/Tests.jspa#/testCase/LVV-T2241}{\textit{ LVV-T2241 } }
test case in Jira.

Testing the communications between MTAOS and M1M3, M2, and the hexapods,
via the MTAOS addAberration command.\\
This test exercises the following components:

\begin{itemize}
\tightlist
\item
  M1M3
\item
  M2
\item
  camera hexapod
\item
  M2 hexapod
\item
  MTAOS
\end{itemize}

This tests that adding aberrations in the way of 1um+1um =
2um.\\[2\baselineskip]Same requirements in
\href{https://jira.lsstcorp.org/secure/Tests.jspa\#/testCase/LVV-T2190}{LVV-2190}.
Different ways to verify for them.\\[2\baselineskip]Requirements are
from \citeds{LTS-186} (MTAOS), \citeds{LTS-206}, \citeds{LTS-88}, \citeds{LTS-146}\\
(MTAOS calculates and publishes the offsets.)\\
MTAOS commands the offsets to the other components\\
LTS- 160,161,162 (Components subscribe to the MTAOS telemetry and apply
the offsets.)

\textbf{ Preconditions}:\\
\begin{itemize}
\tightlist
\item
  M1M3, M2, and Camera Hexapod pass minimum functionality test (we use
  M2 hexapod simulator)
\item
  summit network and EFD are ready to record Camera Hexapods commands,
  events, and telemetry (use of domain 0)
\item
  Jupyter notebook for the test has been written and tested in the NCSA
  simulated environment
\item
  MTAOS CSC is set up ready on the summit network
\end{itemize}

\paragraph{ LVV-T2214 - MTMount Elevation Changes with MTAOS Aberrations }\mbox{}\\

Version \textbf{1}.
Open  \href{https://jira.lsstcorp.org/secure/Tests.jspa#/testCase/LVV-T2214}{\textit{ LVV-T2214 } }
test case in Jira.



\textbf{ Preconditions}:\\


\paragraph{ LVV-T2215 - IM(D): Integrated Slew Test }\mbox{}\\

Version \textbf{2}.
Open  \href{https://jira.lsstcorp.org/secure/Tests.jspa#/testCase/LVV-T2215}{\textit{ LVV-T2215 } }
test case in Jira.

In this test case, we are slewing to four different targets
\textbf{without tracking}.\\
The success criterium is to slew to the targets.\\[2\baselineskip]For
this test case, we should use the \textbf{hybrid} \textbf{mtmount} (CCW
hardware controlled but the Alt/Az simulated).\\[2\baselineskip]This is
a validation test case, no requirement is verified here.

\textbf{ Preconditions}:\\
\begin{itemize}
\tightlist
\item
  M1M3, M2, and Camera Hexapod pass the minimum functionality test (we
  use the M2 hexapod simulator)
\item
  Jupyter notebook for the test has been written and tested at a test
  stand
\end{itemize}

\paragraph{ LVV-T2216 - Integrated Slew and Tracking Test }\mbox{}\\

Version \textbf{1}.
Open  \href{https://jira.lsstcorp.org/secure/Tests.jspa#/testCase/LVV-T2216}{\textit{ LVV-T2216 } }
test case in Jira.

{This test case is very similar
to~}{\href{https://jira.lsstcorp.org/secure/Tests.jspa\#/testCase/LVV-T2215}{LVV-T2215}}{,
but they are not exactly the same.}\\
{This test involves slew and track while the other one involves only
slew with no tracking.}\\[2\baselineskip]{We first get force values from
M1M3 and M2 using the inclinometers then using the mount telemetry and
compare the results.~}

\textbf{ Preconditions}:\\
The components must have passed the minimum functionality tests.

\subsection{Test Cycle LVV-C173 }

Open test cycle {\it \href{https://jira.lsstcorp.org/secure/Tests.jspa#/testrun/LVV-C173}{Level 3 System Integration: M1M3 Hardware with Simulators}} in Jira.

Test Cycle name: Level 3 System Integration: M1M3 Hardware with Simulators\\
Status: In Progress

Get details from Google doc.\\[2\baselineskip]This uses the MTMount
simulator.

\subsubsection{Software Version/Baseline}
Not provided.

\subsubsection{Configuration}
M2 Mirror stays in the horizontal position. We are using telemetry to
simulate elevation changes.\\

\begin{itemize}
\tightlist
\item
  M1M3 hardware
\item
  M2 simulator
\item
  Camera Hexapod simulator
\item
  Camera Rotator simulator
\item
  M2 Hexapod simulator
\item
  MTMount simulator
\item
  MTAOS software
\item
  {ComCam hardware but not utilizing {software}}
\end{itemize}

\subsubsection{Test Cases in LVV-C173 Test Cycle}

\paragraph{ LVV-T2190 - MTAOS add aberrations to M1M3+M2+Hexapods }\mbox{}\\

Version \textbf{1}.
Open  \href{https://jira.lsstcorp.org/secure/Tests.jspa#/testCase/LVV-T2190}{\textit{ LVV-T2190 } }
test case in Jira.

Testing the communications between MTAOS and M1M3, M2, and the hexapods,
via the MTAOS \emph{addAberration~}command.\\
This test exercises the following components:

\begin{itemize}
\tightlist
\item
  M1M3
\item
  M2
\item
  Camera Hexapod
\item
  M2 Hexapod
\item
  MTAOS
\end{itemize}

These test steps are to test adding aberrations in the way of 1um*2um =
2um.\\
Same requirements as in
\href{https://jira.lsstcorp.org/secure/Tests.jspa\#/testCase/LVV-T2241}{LVV-2241}.
Different ways to verify for them.\\[2\baselineskip]Requirements are
from \citeds{LTS-186} (MTAOS), \citeds{LTS-206}, \citeds{LTS-88}, and \citeds{LTS-146} (MTAOS calculates and
publishes the offsets.)\\
MTAOS commands the offsets to the other components LTS- 160,161,162
(Components subscribe to the MTAOS telemetry and apply the offsets.)

\textbf{ Preconditions}:\\
\begin{itemize}
\tightlist
\item
  M1M3, M2, and Camera Hexapod pass the minimum functionality test (we
  use the M2 hexapod simulator)
\item
  Summit network and EFD are ready to record Camera Hexapod's commands,
  events, and telemetry (use of domain 0)
\item
  Jupyter notebook for the test has been written and tested a test stand
\item
  MTAOS CSC is set up and ready on the summit network
\end{itemize}

\paragraph{ LVV-T2193 - MTAOS handling of rejected commands }\mbox{}\\

Version \textbf{1}.
Open  \href{https://jira.lsstcorp.org/secure/Tests.jspa#/testCase/LVV-T2193}{\textit{ LVV-T2193 } }
test case in Jira.

Testing the communications between MTAOS and M1M3, M2, and the hexapods,
when one set of MTAOS corrections gets rejected.\\
This test exercises the following components:

\begin{itemize}
\tightlist
\item
  M1M3
\item
  M2
\item
  Camera Hexapod
\item
  M2 Hexapod
\end{itemize}

Extention of LVV-T2190. It has the same validation ticket traced to it.

\textbf{ Preconditions}:\\
\begin{itemize}
\tightlist
\item
  M1M3, M2, and Camera Hexapod pass the minimum functionality tests (we
  use M2 hexapod simulator as long as necessary)
\item
  Summit network and EFD are ready to record Camera Hexapods commands,
  events, and telemetry (use of domain 0)
\item
  Jupyter notebook for the test has been written and tested in the NCSA
  simulated environment
\item
  MTAOS CSC is set up and ready on the summit network
\end{itemize}

\paragraph{ LVV-T2213 - Look-up Table Application from MTMount Elevation and Azimuth Changes }\mbox{}\\

Version \textbf{1}.
Open  \href{https://jira.lsstcorp.org/secure/Tests.jspa#/testCase/LVV-T2213}{\textit{ LVV-T2213 } }
test case in Jira.

Demonstrate that M1M3, M2, and the hexapods can subscribe to MTMount
elevation changes, and the LUT is applied according to the MTMount
elevation angle.\\
Testing that the components are subscribing to the mount
telemetry\\[2\baselineskip]\citeds{LTS-88} for M1M3\\
\citeds{LTS-146} for M2\\
\citeds{LTS-206} for the hexapods\\
XML for the Mount\\
\citeds{LTS-160}\\
\citeds{LTS-161}\\
\citeds{LTS-162}

\textbf{ Preconditions}:\\
The following components are ready for testing (depending on which test
cycle this is being executed in, each component is either a hardware
component or a simulator):\\

\begin{itemize}
\tightlist
\item
  M1M3
\item
  M2~
\item
  M2 Hexapod
\item
  Camera Hexapod
\item
  MTMount
\end{itemize}

\paragraph{ LVV-T2214 - MTMount Elevation Changes with MTAOS Aberrations }\mbox{}\\

Version \textbf{1}.
Open  \href{https://jira.lsstcorp.org/secure/Tests.jspa#/testCase/LVV-T2214}{\textit{ LVV-T2214 } }
test case in Jira.



\textbf{ Preconditions}:\\


\paragraph{ LVV-T2215 - IM(D): Integrated Slew Test }\mbox{}\\

Version \textbf{2}.
Open  \href{https://jira.lsstcorp.org/secure/Tests.jspa#/testCase/LVV-T2215}{\textit{ LVV-T2215 } }
test case in Jira.

In this test case, we are slewing to four different targets
\textbf{without tracking}.\\
The success criterium is to slew to the targets.\\[2\baselineskip]For
this test case, we should use the \textbf{hybrid} \textbf{mtmount} (CCW
hardware controlled but the Alt/Az simulated).\\[2\baselineskip]This is
a validation test case, no requirement is verified here.

\textbf{ Preconditions}:\\
\begin{itemize}
\tightlist
\item
  M1M3, M2, and Camera Hexapod pass the minimum functionality test (we
  use the M2 hexapod simulator)
\item
  Jupyter notebook for the test has been written and tested at a test
  stand
\end{itemize}

\paragraph{ LVV-T2216 - Integrated Slew and Tracking Test }\mbox{}\\

Version \textbf{1}.
Open  \href{https://jira.lsstcorp.org/secure/Tests.jspa#/testCase/LVV-T2216}{\textit{ LVV-T2216 } }
test case in Jira.

{This test case is very similar
to~}{\href{https://jira.lsstcorp.org/secure/Tests.jspa\#/testCase/LVV-T2215}{LVV-T2215}}{,
but they are not exactly the same.}\\
{This test involves slew and track while the other one involves only
slew with no tracking.}\\[2\baselineskip]{We first get force values from
M1M3 and M2 using the inclinometers then using the mount telemetry and
compare the results.~}

\textbf{ Preconditions}:\\
The components must have passed the minimum functionality tests.

\paragraph{ LVV-T2241 - MTAOS corrections accumulation }\mbox{}\\

Version \textbf{1}.
Open  \href{https://jira.lsstcorp.org/secure/Tests.jspa#/testCase/LVV-T2241}{\textit{ LVV-T2241 } }
test case in Jira.

Testing the communications between MTAOS and M1M3, M2, and the hexapods,
via the MTAOS addAberration command.\\
This test exercises the following components:

\begin{itemize}
\tightlist
\item
  M1M3
\item
  M2
\item
  camera hexapod
\item
  M2 hexapod
\item
  MTAOS
\end{itemize}

This tests that adding aberrations in the way of 1um+1um =
2um.\\[2\baselineskip]Same requirements in
\href{https://jira.lsstcorp.org/secure/Tests.jspa\#/testCase/LVV-T2190}{LVV-2190}.
Different ways to verify for them.\\[2\baselineskip]Requirements are
from \citeds{LTS-186} (MTAOS), \citeds{LTS-206}, \citeds{LTS-88}, \citeds{LTS-146}\\
(MTAOS calculates and publishes the offsets.)\\
MTAOS commands the offsets to the other components\\
LTS- 160,161,162 (Components subscribe to the MTAOS telemetry and apply
the offsets.)

\textbf{ Preconditions}:\\
\begin{itemize}
\tightlist
\item
  M1M3, M2, and Camera Hexapod pass minimum functionality test (we use
  M2 hexapod simulator)
\item
  summit network and EFD are ready to record Camera Hexapods commands,
  events, and telemetry (use of domain 0)
\item
  Jupyter notebook for the test has been written and tested in the NCSA
  simulated environment
\item
  MTAOS CSC is set up ready on the summit network
\end{itemize}

\paragraph{ LVV-T2242 - MTAOS configuration changes }\mbox{}\\

Version \textbf{1}.
Open  \href{https://jira.lsstcorp.org/secure/Tests.jspa#/testCase/LVV-T2242}{\textit{ LVV-T2242 } }
test case in Jira.

Testing the communications between MTAOS (both WEP and OFC) and M1M3,
M2, and the hexapods, via the MTAOS \emph{addAberration} command.\\
This test exercises the following components:

\begin{itemize}
\tightlist
\item
  M1M3
\item
  M2
\item
  camera hexapod
\item
  M2 hexapod
\item
  MTAOS
\end{itemize}

From
\href{https://docushare.lsst.org/docushare/dsweb/Get/LTS-186}{LTS-186 -
MTAOS: WEP and OFC Requirement Document}:\\
- Kallman filtering parameters\\
- Number of visits\\
- Detector(s) to be used*\\
- Location of the detector*\\
- Source selection parameters ({TBC})*\\
- Basis Set\\
-
\href{https://github.com/lsst-ts/ts_config_mttcs/blob/develop/MTAOS/v2/ofc/comcam/senM_9_19_50.yaml}{Sensitivity
Matrix}*\\
- Reconstructor Matrix\\
- Penalty Matrix\\
- Quality Matrix\\
- Gain\\[2\baselineskip]* These are the configurations that we can
verify on System Spread Integration Tests on Level 3.

\textbf{ Preconditions}:\\
\begin{itemize}
\tightlist
\item
  M1M3, M2, and Camera Hexapod pass the minimum functionality tests (we
  use the M2 hexapod simulator)
\item
  Summit network and EFD are ready to record CamHex commands, events,
  and telemetry (use of domain 0)
\item
  Jupyter notebook for the test has been written and tested in the test
  stand
\item
  MTAOS CSC is set up and ready on the summit network
\end{itemize}

\subsection{Test Cycle LVV-C174 }

Open test cycle {\it \href{https://jira.lsstcorp.org/secure/Tests.jspa#/testrun/LVV-C174}{Level 3 System Integration: Camera Hexapod Hardware with M2 Hardware,
M1M3 Hardware, and Simulators}} in Jira.

Test Cycle name: Level 3 System Integration: Camera Hexapod Hardware with M2 Hardware,
M1M3 Hardware, and Simulators\\
Status: In Progress

Get details from Google doc.\\[3\baselineskip]This uses the MTMount
simulator.

\subsubsection{Software Version/Baseline}
Not provided.

\subsubsection{Configuration}
\begin{itemize}
\tightlist
\item
  GIS hardware
\item
  M1M3 hardware
\item
  M2 hardware
\item
  Camera Hexapod hardware
\item
  Camera Rotator simulator
\item
  M2 Hexapod simulator
\item
  MTMount simulator
\item
  MTAOS software
\item
  ComCam hardware but not utilizing software
\end{itemize}

\subsubsection{Test Cases in LVV-C174 Test Cycle}

\paragraph{ LVV-T2193 - MTAOS handling of rejected commands }\mbox{}\\

Version \textbf{1}.
Open  \href{https://jira.lsstcorp.org/secure/Tests.jspa#/testCase/LVV-T2193}{\textit{ LVV-T2193 } }
test case in Jira.

Testing the communications between MTAOS and M1M3, M2, and the hexapods,
when one set of MTAOS corrections gets rejected.\\
This test exercises the following components:

\begin{itemize}
\tightlist
\item
  M1M3
\item
  M2
\item
  Camera Hexapod
\item
  M2 Hexapod
\end{itemize}

Extention of LVV-T2190. It has the same validation ticket traced to it.

\textbf{ Preconditions}:\\
\begin{itemize}
\tightlist
\item
  M1M3, M2, and Camera Hexapod pass the minimum functionality tests (we
  use M2 hexapod simulator as long as necessary)
\item
  Summit network and EFD are ready to record Camera Hexapods commands,
  events, and telemetry (use of domain 0)
\item
  Jupyter notebook for the test has been written and tested in the NCSA
  simulated environment
\item
  MTAOS CSC is set up and ready on the summit network
\end{itemize}

\paragraph{ LVV-T2242 - MTAOS configuration changes }\mbox{}\\

Version \textbf{1}.
Open  \href{https://jira.lsstcorp.org/secure/Tests.jspa#/testCase/LVV-T2242}{\textit{ LVV-T2242 } }
test case in Jira.

Testing the communications between MTAOS (both WEP and OFC) and M1M3,
M2, and the hexapods, via the MTAOS \emph{addAberration} command.\\
This test exercises the following components:

\begin{itemize}
\tightlist
\item
  M1M3
\item
  M2
\item
  camera hexapod
\item
  M2 hexapod
\item
  MTAOS
\end{itemize}

From
\href{https://docushare.lsst.org/docushare/dsweb/Get/LTS-186}{LTS-186 -
MTAOS: WEP and OFC Requirement Document}:\\
- Kallman filtering parameters\\
- Number of visits\\
- Detector(s) to be used*\\
- Location of the detector*\\
- Source selection parameters ({TBC})*\\
- Basis Set\\
-
\href{https://github.com/lsst-ts/ts_config_mttcs/blob/develop/MTAOS/v2/ofc/comcam/senM_9_19_50.yaml}{Sensitivity
Matrix}*\\
- Reconstructor Matrix\\
- Penalty Matrix\\
- Quality Matrix\\
- Gain\\[2\baselineskip]* These are the configurations that we can
verify on System Spread Integration Tests on Level 3.

\textbf{ Preconditions}:\\
\begin{itemize}
\tightlist
\item
  M1M3, M2, and Camera Hexapod pass the minimum functionality tests (we
  use the M2 hexapod simulator)
\item
  Summit network and EFD are ready to record CamHex commands, events,
  and telemetry (use of domain 0)
\item
  Jupyter notebook for the test has been written and tested in the test
  stand
\item
  MTAOS CSC is set up and ready on the summit network
\end{itemize}

\paragraph{ LVV-T2213 - Look-up Table Application from MTMount Elevation and Azimuth Changes }\mbox{}\\

Version \textbf{1}.
Open  \href{https://jira.lsstcorp.org/secure/Tests.jspa#/testCase/LVV-T2213}{\textit{ LVV-T2213 } }
test case in Jira.

Demonstrate that M1M3, M2, and the hexapods can subscribe to MTMount
elevation changes, and the LUT is applied according to the MTMount
elevation angle.\\
Testing that the components are subscribing to the mount
telemetry\\[2\baselineskip]\citeds{LTS-88} for M1M3\\
\citeds{LTS-146} for M2\\
\citeds{LTS-206} for the hexapods\\
XML for the Mount\\
\citeds{LTS-160}\\
\citeds{LTS-161}\\
\citeds{LTS-162}

\textbf{ Preconditions}:\\
The following components are ready for testing (depending on which test
cycle this is being executed in, each component is either a hardware
component or a simulator):\\

\begin{itemize}
\tightlist
\item
  M1M3
\item
  M2~
\item
  M2 Hexapod
\item
  Camera Hexapod
\item
  MTMount
\end{itemize}

\paragraph{ LVV-T2190 - MTAOS add aberrations to M1M3+M2+Hexapods }\mbox{}\\

Version \textbf{1}.
Open  \href{https://jira.lsstcorp.org/secure/Tests.jspa#/testCase/LVV-T2190}{\textit{ LVV-T2190 } }
test case in Jira.

Testing the communications between MTAOS and M1M3, M2, and the hexapods,
via the MTAOS \emph{addAberration~}command.\\
This test exercises the following components:

\begin{itemize}
\tightlist
\item
  M1M3
\item
  M2
\item
  Camera Hexapod
\item
  M2 Hexapod
\item
  MTAOS
\end{itemize}

These test steps are to test adding aberrations in the way of 1um*2um =
2um.\\
Same requirements as in
\href{https://jira.lsstcorp.org/secure/Tests.jspa\#/testCase/LVV-T2241}{LVV-2241}.
Different ways to verify for them.\\[2\baselineskip]Requirements are
from \citeds{LTS-186} (MTAOS), \citeds{LTS-206}, \citeds{LTS-88}, and \citeds{LTS-146} (MTAOS calculates and
publishes the offsets.)\\
MTAOS commands the offsets to the other components LTS- 160,161,162
(Components subscribe to the MTAOS telemetry and apply the offsets.)

\textbf{ Preconditions}:\\
\begin{itemize}
\tightlist
\item
  M1M3, M2, and Camera Hexapod pass the minimum functionality test (we
  use the M2 hexapod simulator)
\item
  Summit network and EFD are ready to record Camera Hexapod's commands,
  events, and telemetry (use of domain 0)
\item
  Jupyter notebook for the test has been written and tested a test stand
\item
  MTAOS CSC is set up and ready on the summit network
\end{itemize}

\paragraph{ LVV-T2241 - MTAOS corrections accumulation }\mbox{}\\

Version \textbf{1}.
Open  \href{https://jira.lsstcorp.org/secure/Tests.jspa#/testCase/LVV-T2241}{\textit{ LVV-T2241 } }
test case in Jira.

Testing the communications between MTAOS and M1M3, M2, and the hexapods,
via the MTAOS addAberration command.\\
This test exercises the following components:

\begin{itemize}
\tightlist
\item
  M1M3
\item
  M2
\item
  camera hexapod
\item
  M2 hexapod
\item
  MTAOS
\end{itemize}

This tests that adding aberrations in the way of 1um+1um =
2um.\\[2\baselineskip]Same requirements in
\href{https://jira.lsstcorp.org/secure/Tests.jspa\#/testCase/LVV-T2190}{LVV-2190}.
Different ways to verify for them.\\[2\baselineskip]Requirements are
from \citeds{LTS-186} (MTAOS), \citeds{LTS-206}, \citeds{LTS-88}, \citeds{LTS-146}\\
(MTAOS calculates and publishes the offsets.)\\
MTAOS commands the offsets to the other components\\
LTS- 160,161,162 (Components subscribe to the MTAOS telemetry and apply
the offsets.)

\textbf{ Preconditions}:\\
\begin{itemize}
\tightlist
\item
  M1M3, M2, and Camera Hexapod pass minimum functionality test (we use
  M2 hexapod simulator)
\item
  summit network and EFD are ready to record Camera Hexapods commands,
  events, and telemetry (use of domain 0)
\item
  Jupyter notebook for the test has been written and tested in the NCSA
  simulated environment
\item
  MTAOS CSC is set up ready on the summit network
\end{itemize}

\paragraph{ LVV-T2214 - MTMount Elevation Changes with MTAOS Aberrations }\mbox{}\\

Version \textbf{1}.
Open  \href{https://jira.lsstcorp.org/secure/Tests.jspa#/testCase/LVV-T2214}{\textit{ LVV-T2214 } }
test case in Jira.



\textbf{ Preconditions}:\\


\paragraph{ LVV-T2215 - IM(D): Integrated Slew Test }\mbox{}\\

Version \textbf{2}.
Open  \href{https://jira.lsstcorp.org/secure/Tests.jspa#/testCase/LVV-T2215}{\textit{ LVV-T2215 } }
test case in Jira.

In this test case, we are slewing to four different targets
\textbf{without tracking}.\\
The success criterium is to slew to the targets.\\[2\baselineskip]For
this test case, we should use the \textbf{hybrid} \textbf{mtmount} (CCW
hardware controlled but the Alt/Az simulated).\\[2\baselineskip]This is
a validation test case, no requirement is verified here.

\textbf{ Preconditions}:\\
\begin{itemize}
\tightlist
\item
  M1M3, M2, and Camera Hexapod pass the minimum functionality test (we
  use the M2 hexapod simulator)
\item
  Jupyter notebook for the test has been written and tested at a test
  stand
\end{itemize}

\paragraph{ LVV-T2216 - Integrated Slew and Tracking Test }\mbox{}\\

Version \textbf{1}.
Open  \href{https://jira.lsstcorp.org/secure/Tests.jspa#/testCase/LVV-T2216}{\textit{ LVV-T2216 } }
test case in Jira.

{This test case is very similar
to~}{\href{https://jira.lsstcorp.org/secure/Tests.jspa\#/testCase/LVV-T2215}{LVV-T2215}}{,
but they are not exactly the same.}\\
{This test involves slew and track while the other one involves only
slew with no tracking.}\\[2\baselineskip]{We first get force values from
M1M3 and M2 using the inclinometers then using the mount telemetry and
compare the results.~}

\textbf{ Preconditions}:\\
The components must have passed the minimum functionality tests.

\subsection{Test Cycle LVV-C175 }

Open test cycle {\it \href{https://jira.lsstcorp.org/secure/Tests.jspa#/testrun/LVV-C175}{Level 3 System Integration: Camera Rotator Hardware with Camera Hexapod
Hardware, M2 Hardware, M1M3 Hardware, and Simulators}} in Jira.

Test Cycle name: Level 3 System Integration: Camera Rotator Hardware with Camera Hexapod
Hardware, M2 Hardware, M1M3 Hardware, and Simulators\\
Status: In Progress

Get details from Google doc.\\[2\baselineskip]This test is using the
MTMount \textbf{simulated.}

\subsubsection{Software Version/Baseline}
Not provided.

\subsubsection{Configuration}
\begin{itemize}
\tightlist
\item
  GIS hardware
\item
  M1M3 hardware
\item
  M2 hardware
\item
  Camera Hexapod hardware
\item
  Camera Rotator hardware
\item
  M2 Hexapod simulator
\item
  MTMount simulator
\item
  MTAOS software
\item
  ComCam hardware but not utilizing software
\end{itemize}

\subsubsection{Test Cases in LVV-C175 Test Cycle}

\paragraph{ LVV-T2215 - IM(D): Integrated Slew Test }\mbox{}\\

Version \textbf{2}.
Open  \href{https://jira.lsstcorp.org/secure/Tests.jspa#/testCase/LVV-T2215}{\textit{ LVV-T2215 } }
test case in Jira.

In this test case, we are slewing to four different targets
\textbf{without tracking}.\\
The success criterium is to slew to the targets.\\[2\baselineskip]For
this test case, we should use the \textbf{hybrid} \textbf{mtmount} (CCW
hardware controlled but the Alt/Az simulated).\\[2\baselineskip]This is
a validation test case, no requirement is verified here.

\textbf{ Preconditions}:\\
\begin{itemize}
\tightlist
\item
  M1M3, M2, and Camera Hexapod pass the minimum functionality test (we
  use the M2 hexapod simulator)
\item
  Jupyter notebook for the test has been written and tested at a test
  stand
\end{itemize}

\paragraph{ LVV-T2216 - Integrated Slew and Tracking Test }\mbox{}\\

Version \textbf{1}.
Open  \href{https://jira.lsstcorp.org/secure/Tests.jspa#/testCase/LVV-T2216}{\textit{ LVV-T2216 } }
test case in Jira.

{This test case is very similar
to~}{\href{https://jira.lsstcorp.org/secure/Tests.jspa\#/testCase/LVV-T2215}{LVV-T2215}}{,
but they are not exactly the same.}\\
{This test involves slew and track while the other one involves only
slew with no tracking.}\\[2\baselineskip]{We first get force values from
M1M3 and M2 using the inclinometers then using the mount telemetry and
compare the results.~}

\textbf{ Preconditions}:\\
The components must have passed the minimum functionality tests.

\subsection{Test Cycle LVV-C176 }

Open test cycle {\it \href{https://jira.lsstcorp.org/secure/Tests.jspa#/testrun/LVV-C176}{Level 3 System Integration: ComCam Software with Camera Rotator
Hardware, Camera Hexapod Hardware, M2 Hardware, M1M3 Hardware, and
Simulators}} in Jira.

Test Cycle name: Level 3 System Integration: ComCam Software with Camera Rotator
Hardware, Camera Hexapod Hardware, M2 Hardware, M1M3 Hardware, and
Simulators\\
Status: In Progress

The test case aims to verify that the components\\

\begin{itemize}
\tightlist
\item
  GIS hardware
\item
  M1M3 hardware
\item
  M2 hardware
\item
  Camera Hexapod hardware
\item
  Camera Rotator hardware
\item
  M2 Hexapod simulator
\item
  MTMount simulator with CCW hardware
\end{itemize}

are seamlessly integrated with the\\

\begin{itemize}
\tightlist
\item
  MTAOS software
\item
  Pointing Component software
\item
  ComCam hardware and software
\end{itemize}

This test cycle needs to use the MTMount \textbf{hybrid.}

\subsubsection{Software Version/Baseline}
Not provided.

\subsubsection{Configuration}
\begin{itemize}
\tightlist
\item
  GIS hardware
\item
  M1M3 hardware
\item
  M2 hardware
\item
  Camera Hexapod hardware
\item
  Camera Rotator hardware
\item
  M2 Hexapod simulator
\item
  MTMount simulator with CCW hardware
\item
  MTAOS software
\item
  Pointing Component software
\item
  ComCam hardware and software
\end{itemize}

\subsubsection{Test Cases in LVV-C176 Test Cycle}

\paragraph{ LVV-T2578 - ComCam Opto-Mechanical Filter Changer Test }\mbox{}\\

Version \textbf{1}.
Open  \href{https://jira.lsstcorp.org/secure/Tests.jspa#/testCase/LVV-T2578}{\textit{ LVV-T2578 } }
test case in Jira.

The objective of this test is to verify the \citeds{LTS-508} requirements in
relation to the filter changer performance.

\textbf{ Preconditions}:\\
ComCam operable.

\paragraph{ LVV-T2344 - Startup MT Components for System Spread Integration Tests on Level 3 via
Jupyter Notebook }\mbox{}\\

Version \textbf{1}.
Open  \href{https://jira.lsstcorp.org/secure/Tests.jspa#/testCase/LVV-T2344}{\textit{ LVV-T2344 } }
test case in Jira.

Concentrate the setup procedures for all the components used during
Phases 1, 2, and 3 of the Level 3 spread configuration integration
tests.\\
Based on the no-longer existing~\textbf{setup.ipynb} notebook (it was
updated and renamed, see the Jupyter Notebook Link further below).

\textbf{ Preconditions}:\\


\paragraph{ LVV-T2190 - MTAOS add aberrations to M1M3+M2+Hexapods }\mbox{}\\

Version \textbf{1}.
Open  \href{https://jira.lsstcorp.org/secure/Tests.jspa#/testCase/LVV-T2190}{\textit{ LVV-T2190 } }
test case in Jira.

Testing the communications between MTAOS and M1M3, M2, and the hexapods,
via the MTAOS \emph{addAberration~}command.\\
This test exercises the following components:

\begin{itemize}
\tightlist
\item
  M1M3
\item
  M2
\item
  Camera Hexapod
\item
  M2 Hexapod
\item
  MTAOS
\end{itemize}

These test steps are to test adding aberrations in the way of 1um*2um =
2um.\\
Same requirements as in
\href{https://jira.lsstcorp.org/secure/Tests.jspa\#/testCase/LVV-T2241}{LVV-2241}.
Different ways to verify for them.\\[2\baselineskip]Requirements are
from \citeds{LTS-186} (MTAOS), \citeds{LTS-206}, \citeds{LTS-88}, and \citeds{LTS-146} (MTAOS calculates and
publishes the offsets.)\\
MTAOS commands the offsets to the other components LTS- 160,161,162
(Components subscribe to the MTAOS telemetry and apply the offsets.)

\textbf{ Preconditions}:\\
\begin{itemize}
\tightlist
\item
  M1M3, M2, and Camera Hexapod pass the minimum functionality test (we
  use the M2 hexapod simulator)
\item
  Summit network and EFD are ready to record Camera Hexapod's commands,
  events, and telemetry (use of domain 0)
\item
  Jupyter notebook for the test has been written and tested a test stand
\item
  MTAOS CSC is set up and ready on the summit network
\end{itemize}

\paragraph{ LVV-T2290 - Slew, Track and Image taking with ComCam }\mbox{}\\

Version \textbf{2}.
Open  \href{https://jira.lsstcorp.org/secure/Tests.jspa#/testCase/LVV-T2290}{\textit{ LVV-T2290 } }
test case in Jira.

Use the procedures from the TCS
(https://github.com/lsst-ts/ts\_notebooks/tree/tickets/DM-31475/procedures/MT/slew)
to set up, slew, track and take an image with ComCam\\
The success criteria are to obtain an image with ComCam, available in
the Butler and populated with the correct metadata.\\
\citeds{LSE-71} OCS to Camera
(\href{https://jira.lsstcorp.org/browse/LVV-5129}{LVV-5129~}+4)\\
Requirements for the Butler\\[2\baselineskip]Find data by searching for
header information (metadata information) in the Butler.
--\textgreater{} make a new test case.\\[2\baselineskip]Access image
from the butler.\\
Metadata is available and has the available data at the moment of the
test. Date, Type of image, mount parameters(EL, AZ), Pointing(RA, DEC)

\textbf{ Preconditions}:\\
\begin{itemize}
\tightlist
\item
  M1M3, M2, and Camera Hexapod pass the minimum functionality test (we
  use the M2 hexapod simulator)
\item
  Summit network and EFD are ready to record Camera Hexapod's commands,
  events, and telemetry (use of domain 0)
\item
  Jupyter notebook for the test has been written and tested at a test
  stand.
\end{itemize}

\paragraph{ LVV-T2193 - MTAOS handling of rejected commands }\mbox{}\\

Version \textbf{1}.
Open  \href{https://jira.lsstcorp.org/secure/Tests.jspa#/testCase/LVV-T2193}{\textit{ LVV-T2193 } }
test case in Jira.

Testing the communications between MTAOS and M1M3, M2, and the hexapods,
when one set of MTAOS corrections gets rejected.\\
This test exercises the following components:

\begin{itemize}
\tightlist
\item
  M1M3
\item
  M2
\item
  Camera Hexapod
\item
  M2 Hexapod
\end{itemize}

Extention of LVV-T2190. It has the same validation ticket traced to it.

\textbf{ Preconditions}:\\
\begin{itemize}
\tightlist
\item
  M1M3, M2, and Camera Hexapod pass the minimum functionality tests (we
  use M2 hexapod simulator as long as necessary)
\item
  Summit network and EFD are ready to record Camera Hexapods commands,
  events, and telemetry (use of domain 0)
\item
  Jupyter notebook for the test has been written and tested in the NCSA
  simulated environment
\item
  MTAOS CSC is set up and ready on the summit network
\end{itemize}

\paragraph{ LVV-T2228 - IM(G): One time ComCam Image Ingestion and MTAOS Correction }\mbox{}\\

Version \textbf{1}.
Open  \href{https://jira.lsstcorp.org/secure/Tests.jspa#/testCase/LVV-T2228}{\textit{ LVV-T2228 } }
test case in Jira.

The goal of this test case is to ingest a simulated image using the
playback mode of ComCam.\\
This means that even though simulated, the system after the DAQ behaves
as if it were real.\\
This includes the DM pipeline (Butler, OCPS, etc) and, in a later stage,
the AOS correction loop.\\
In this particular test case, we are focusing on the Active Optics
pipeline.\\[2\baselineskip]{Demonstrate:}\\
{•}{~}{Slew to an observing field, generating telemetry
(}{cf.}{~}{IM}{F}{)}\\
{•}{~}{piston the camera in Z-direction (e.g. in, -out, in, +out,
in).}\\
{•}{~}{take data using ComCam in playback mode, commanded from the
OCS}\\
{•}{~}{uses}{~}{OCPS}{~}{to reduce the data on the commissioning
cluster}\\
{•}{~}{load the Zernike coefficients into}{~the~}{EFD}\\
{•}{~}{Use the MTAOS to update the state of the
optics}\\[2\baselineskip]{Associated Documents}

\begin{itemize}
\tightlist
\item
  \href{https://ls.st/Document-14771}{Document-14771 - Active Optics
  Baseline Design}
\item
  \href{https://ls.st/LTS-186}{LSE-186 - Main Telescope Active Optics
  System: Wavefront Estimation Pipeline and Optical Feedback Controller
  Requirement Document}
\end{itemize}

\textbf{ Preconditions}:\\
IM1 has been successfully completed:
\href{https://jira.lsstcorp.org/browse/SUMMIT-2983}{SUMMIT-2983}\\
The precondition is fulfilled since Nov 2020.

\paragraph{ LVV-T2229 - Closed Loop ComCam Image Ingestion and Application of Correction }\mbox{}\\

Version \textbf{2}.
Open  \href{https://jira.lsstcorp.org/secure/Tests.jspa#/testCase/LVV-T2229}{\textit{ LVV-T2229 } }
test case in Jira.

The goal of this test case is to\\[2\baselineskip]

\begin{itemize}
\tightlist
\item
  perform three telescope slews including the ComCam software and
  hardware.
\item
  apply the calculated offsets to the different AOS components from the
  first visit.
\end{itemize}

This means that even though simulated, the system after the DAQ behaves
as if it were real. This includes the DM pipeline (Butler, OCPS, etc)
and then further down the AOS loop.~\\
This should follow after step ``15 IMG {[}g{]}: AOS using ComCam,
commanding MT (2021-06-30)'' in~\url{https://sitcomtn-006.lsst.io/},
revision: 2021-09-03.

\textbf{ Preconditions}:\\
One precondition for this test is to have the DAQ playback capability.

\subsection{Test Cycle LVV-C177 }

Open test cycle {\it \href{https://jira.lsstcorp.org/secure/Tests.jspa#/testrun/LVV-C177}{Level 3 System Integration: GIS Integration}} in Jira.

Test Cycle name: Level 3 System Integration: GIS Integration\\
Status: Not Executed

Verify that all the GIS interlocks for the available hardware and
simulated inputs/outputs for the unavailable hardware are fully
functional.

\subsubsection{Software Version/Baseline}
Not provided.

\subsubsection{Configuration}
No changing configuration for the test cases within this test cycle.

\subsubsection{Test Cases in LVV-C177 Test Cycle}

\paragraph{ LVV-T2540 - Level 3 Interlock Systems Inspections }\mbox{}\\

Version \textbf{1}.
Open  \href{https://jira.lsstcorp.org/secure/Tests.jspa#/testCase/LVV-T2540}{\textit{ LVV-T2540 } }
test case in Jira.

Inspect the following Interlock Systems that are located on Level 3:

\begin{itemize}
\tightlist
\item
  Camera Hexapod
\item
  Camera Rotator
\item
  M2
\item
  M1M3
\item
  GIS
\end{itemize}

\textbf{ Preconditions}:\\
All of the hardware for the Interlock Systems are located on Level 3 and
connected through the GIS.

\paragraph{ LVV-T2227 - GIS Safety Matrix Interlock Verification }\mbox{}\\

Version \textbf{1}.
Open  \href{https://jira.lsstcorp.org/secure/Tests.jspa#/testCase/LVV-T2227}{\textit{ LVV-T2227 } }
test case in Jira.

The objective of this test case is to trigger each detection identified
in the GIS Safety Matrix per \citeds{LTS-188} and verify that the corresponding
actuations are generated. Depending on the detections, actuations, and
hardware available on the third floor of the Summit some of the
detections will be triggered by simulating the input to the GIS hardware
and actuation outputs from the GIS hardware will be verified using test
equipment in lieu of the actual hardware. If the Interlock System (IS)
for a system that interfaces with the GIS is available on the third
floor of the Summit, the input from the GIS will be verified to be
received by the IS and the corresponding actuations for that IS will be
verified to have been generated and the hardware responds accordingly.

\textbf{ Preconditions}:\\


\paragraph{ LVV-T2231 - (GIS-D-3) GIS Fault Verification }\mbox{}\\

Version \textbf{1}.
Open  \href{https://jira.lsstcorp.org/secure/Tests.jspa#/testCase/LVV-T2231}{\textit{ LVV-T2231 } }
test case in Jira.

The objective of this test case is to trigger each GIS fault identified
in \citeds{LTS-188} and verify that the corresponding actuations are generated.
Depending on the detections, actuations, and hardware available on the
third floor of the Summit some of the detections will be triggered by
simulating the input to the GIS hardware and actuation outputs from the
GIS hardware will be verified using test equipment in lieu of the actual
hardware. If the Interlock System (IS) for a system that interfaces with
the GIS is available on the third floor of the Summit, the input from
the GIS will be verified to be received by the IS and the corresponding
actuations for that IS will be verified to have been generated and the
hardware responds accordingly.

\textbf{ Preconditions}:\\


\subsection{Test Cycle LVV-C206 }

Open test cycle {\it \href{https://jira.lsstcorp.org/secure/Tests.jspa#/testrun/LVV-C206}{Level 3 System Integration: M2 Hardware, M1M3 Hardware, CCW+Rotator
Hardware, Cam Hex and Simulators}} in Jira.

Test Cycle name: Level 3 System Integration: M2 Hardware, M1M3 Hardware, CCW+Rotator
Hardware, Cam Hex and Simulators\\
Status: In Progress

In June 2022, during the System spread testing, we tried to use the
simulators to conduct LVV-C173, LVV-C174 and LVV-C175 but the Mount
simulator was producing erroneous elevation data causing the M2 hardware
and M1M3 hardware to fault.\\[2\baselineskip]This test cycle is only to
be performed if the mount simulator can not be debugged on time for the
execution and the M2 simulator can not be deployed.\\[2\baselineskip]In
this cycle, we have tested that we can slew and track while applying the
LUT to the mirror system (m1m3 and m2 subscribed to the mount instead of
the inclinometer).\\[2\baselineskip]This test cycle brings in the CCW
and therefore has to use the MTMount\\[2\baselineskip]Here we are
executing the test cases with larger rotations.\\
TBD: Update the test cases for the conditions of the rotation angle. The
rotation angle should be below 2.5 deg with the hybrid mount and larger
with the hybrid mount.\\[2\baselineskip]

\subsubsection{Software Version/Baseline}
Not provided.

\subsubsection{Configuration}
\begin{itemize}
\tightlist
\item
  GIS hardware
\item
  M1M3 hardware
\item
  M2 hardware
\item
  Camera Hexapod hardware
\item
  Camera Rotator Hardware
\item
  M2 Hexapod simulator
\item
  MTMount simulator with the CCW hardware
\item
  MTAOS software
\item
  ComCam hardware but not utilizing software
\end{itemize}

\subsubsection{Test Cases in LVV-C206 Test Cycle}

\paragraph{ LVV-T2215 - IM(D): Integrated Slew Test }\mbox{}\\

Version \textbf{2}.
Open  \href{https://jira.lsstcorp.org/secure/Tests.jspa#/testCase/LVV-T2215}{\textit{ LVV-T2215 } }
test case in Jira.

In this test case, we are slewing to four different targets
\textbf{without tracking}.\\
The success criterium is to slew to the targets.\\[2\baselineskip]For
this test case, we should use the \textbf{hybrid} \textbf{mtmount} (CCW
hardware controlled but the Alt/Az simulated).\\[2\baselineskip]This is
a validation test case, no requirement is verified here.

\textbf{ Preconditions}:\\
\begin{itemize}
\tightlist
\item
  M1M3, M2, and Camera Hexapod pass the minimum functionality test (we
  use the M2 hexapod simulator)
\item
  Jupyter notebook for the test has been written and tested at a test
  stand
\end{itemize}

\paragraph{ LVV-T2216 - Integrated Slew and Tracking Test }\mbox{}\\

Version \textbf{1}.
Open  \href{https://jira.lsstcorp.org/secure/Tests.jspa#/testCase/LVV-T2216}{\textit{ LVV-T2216 } }
test case in Jira.

{This test case is very similar
to~}{\href{https://jira.lsstcorp.org/secure/Tests.jspa\#/testCase/LVV-T2215}{LVV-T2215}}{,
but they are not exactly the same.}\\
{This test involves slew and track while the other one involves only
slew with no tracking.}\\[2\baselineskip]{We first get force values from
M1M3 and M2 using the inclinometers then using the mount telemetry and
compare the results.~}

\textbf{ Preconditions}:\\
The components must have passed the minimum functionality tests.

\paragraph{ LVV-T2241 - MTAOS corrections accumulation }\mbox{}\\

Version \textbf{1}.
Open  \href{https://jira.lsstcorp.org/secure/Tests.jspa#/testCase/LVV-T2241}{\textit{ LVV-T2241 } }
test case in Jira.

Testing the communications between MTAOS and M1M3, M2, and the hexapods,
via the MTAOS addAberration command.\\
This test exercises the following components:

\begin{itemize}
\tightlist
\item
  M1M3
\item
  M2
\item
  camera hexapod
\item
  M2 hexapod
\item
  MTAOS
\end{itemize}

This tests that adding aberrations in the way of 1um+1um =
2um.\\[2\baselineskip]Same requirements in
\href{https://jira.lsstcorp.org/secure/Tests.jspa\#/testCase/LVV-T2190}{LVV-2190}.
Different ways to verify for them.\\[2\baselineskip]Requirements are
from \citeds{LTS-186} (MTAOS), \citeds{LTS-206}, \citeds{LTS-88}, \citeds{LTS-146}\\
(MTAOS calculates and publishes the offsets.)\\
MTAOS commands the offsets to the other components\\
LTS- 160,161,162 (Components subscribe to the MTAOS telemetry and apply
the offsets.)

\textbf{ Preconditions}:\\
\begin{itemize}
\tightlist
\item
  M1M3, M2, and Camera Hexapod pass minimum functionality test (we use
  M2 hexapod simulator)
\item
  summit network and EFD are ready to record Camera Hexapods commands,
  events, and telemetry (use of domain 0)
\item
  Jupyter notebook for the test has been written and tested in the NCSA
  simulated environment
\item
  MTAOS CSC is set up ready on the summit network
\end{itemize}

\paragraph{ LVV-T2213 - Look-up Table Application from MTMount Elevation and Azimuth Changes }\mbox{}\\

Version \textbf{1}.
Open  \href{https://jira.lsstcorp.org/secure/Tests.jspa#/testCase/LVV-T2213}{\textit{ LVV-T2213 } }
test case in Jira.

Demonstrate that M1M3, M2, and the hexapods can subscribe to MTMount
elevation changes, and the LUT is applied according to the MTMount
elevation angle.\\
Testing that the components are subscribing to the mount
telemetry\\[2\baselineskip]\citeds{LTS-88} for M1M3\\
\citeds{LTS-146} for M2\\
\citeds{LTS-206} for the hexapods\\
XML for the Mount\\
\citeds{LTS-160}\\
\citeds{LTS-161}\\
\citeds{LTS-162}

\textbf{ Preconditions}:\\
The following components are ready for testing (depending on which test
cycle this is being executed in, each component is either a hardware
component or a simulator):\\

\begin{itemize}
\tightlist
\item
  M1M3
\item
  M2~
\item
  M2 Hexapod
\item
  Camera Hexapod
\item
  MTMount
\end{itemize}

\paragraph{ LVV-T2214 - MTMount Elevation Changes with MTAOS Aberrations }\mbox{}\\

Version \textbf{1}.
Open  \href{https://jira.lsstcorp.org/secure/Tests.jspa#/testCase/LVV-T2214}{\textit{ LVV-T2214 } }
test case in Jira.



\textbf{ Preconditions}:\\




\input{appendix.tex}
\end{document}
